\documentclass[11pt]{amsart}
\usepackage[a4paper,margin=1in]{geometry}
\usepackage[T1]{fontenc}
\usepackage{lmodern}
\usepackage{amssymb,amsthm,booktabs}
\usepackage{xurl}
\usepackage[hidelinks]{hyperref}

\newtheorem{theorem}{Theorem}[section]
\newtheorem{lemma}[theorem]{Lemma}
\newtheorem{proposition}[theorem]{Proposition}
\newtheorem{corollary}[theorem]{Corollary}
\newtheorem{fact}[theorem]{Fact}
\newtheorem*{theoremA}{Theorem A}
\newtheorem*{theoremB}{Theorem B}
\theoremstyle{definition}
\newtheorem{definition}[theorem]{Definition}
\newtheorem*{conjecture*}{Conjecture}
\theoremstyle{remark}
\newtheorem{remark}[theorem]{Remark}

\numberwithin{equation}{section}
\setlength{\emergencystretch}{2em}
\exhyphenpenalty=10000 % no line breaks after explicit hyphens (keeps dates and hashes whole)

% ---------------------------------------------------------------------
% Notation
\newcommand{\linv}{\ell_{\mathrm{inv}}}
\newcommand{\IP}{\mathcal{IP}}
\newcommand{\LB}{\mathcal{LB}}
\newcommand{\BR}{\mathcal{BR}}
\newcommand{\LR}{\mathcal{LR}}
\newcommand{\X}{\mathcal{X}}
\newcommand{\leR}{\le_R}
\newcommand{\leL}{\le_L}
\newcommand{\leB}{\le_B}
\newcommand{\geB}{\ge_B}
\newcommand{\abs}[1]{\lvert #1\rvert}
\newcommand{\set}[1]{\{#1\}}
\DeclareMathOperator{\rem}{rem}
\DeclareMathOperator{\disp}{disp}
\DeclareMathOperator{\rank}{rank}
\DeclareMathOperator{\PInv}{PInv}
\DeclareMathOperator{\VInv}{VInv}
\DeclareMathOperator{\card}{card}
\DeclareMathOperator{\invB}{inv}
\DeclareMathOperator{\negB}{neg}
\DeclareMathOperator{\nspB}{nsp}
\DeclareMathOperator{\lb}{lb}
\DeclareMathOperator{\br}{br}
\DeclareMathOperator{\lr}{lr}
\DeclareUrlCommand\file{\urlstyle{tt}}% file paths in this repository; breakable, not hyperlinked

% ---------------------------------------------------------------------
% Zenodo DOI. Leave empty until the DOI is assigned; then insert it here,
% for example \newcommand*{\zenodoDOI}{10.5281/zenodo.1234567}.
% When the macro is empty, nothing is printed.
\newcommand*{\zenodoDOI}{10.5281/zenodo.24613550}
% ---------------------------------------------------------------------
\newcommand*{\doiline}{\ifx\zenodoDOI\empty\else. DOI: \href{https://doi.org/\zenodoDOI}{\nolinkurl{\zenodoDOI}}\fi}

\title[Hart--Rowley Conjecture 1.3 in type A]{Hart--Rowley Conjecture~1.3 holds in type A}
\author{Sukhraj Sachdeva}
\date{2026-11-27\doiline}

\begin{document}

\begin{abstract}
Hart and Rowley conjectured that in a finite Coxeter group every element has involution length at most the rank of the group. The involution length counts the factors of the greedy decomposition of an element into longest involution prefixes. Levinstein has shown that the conjecture is false in general: it fails in the groups \(H_3\), \(F_4\), \(E_6\) and \(E_7\). We prove that it holds in type A. For every permutation \(w\) of \(\{1,\dots,N\}\), the involution length of \(w\) is at most the largest displacement \(\max_i\abs{w(i)-i}\), and this is at most \(N-1\), the rank. The bound is attained. The key step shows that removing the longest involution prefix strictly lowers the displacement. It uses a theorem of Gaetz and Gao on involution prefixes; we write out their proof in full from textbook facts. We also show that for signed permutations the involution length equals the involution length in the symmetric group \(S_{2n}\). This gives the bound \(2n-1\) in type \(B_n\), which is weaker than the rank \(n\). The same reduction fails in type \(D_3\). All proofs are by hand. Exact computations for \(N\le7\) corroborate them. Nothing in this paper is formalized in a proof assistant.
\end{abstract}

\maketitle

\noindent\emph{Disclosure: Human researcher assisted by AI systems during search, proof drafting, and verification. Internally audited and computationally verified, externally unrefereed.}\par\medskip

\section{Introduction}\label{sec:intro}

\subsection{The conjecture}
A Coxeter group is a group generated by a finite set \(S\) of elements of order two, subject only to relations that fix the order of each product of two generators. The number \(\abs{S}\) of generators is the \emph{rank}. Every element \(w\) has a \emph{length} \(\ell(w)\): the least number of generators whose product is \(w\). A \emph{prefix} of \(w\) is an element \(u\) such that some shortest word for \(w\) begins with a shortest word for \(u\). Precise definitions are in Section~\ref{sec:coxeter}.

Hart and Rowley~\cite{HR} studied the prefixes of \(w\) that are involutions. They conjectured that, in every finite Coxeter group, among the involution prefixes of each \(w\neq e\) there is exactly one of largest length, the \emph{ancestor} \(a(w)\)~\cite[Conj.~1.2]{HR}. Write \(w=a(w)\cdot w'\). Repeating the construction on \(w'\), and so on, gives the \emph{ancestor decomposition} of \(w\) into involutions. The number of factors is the \emph{involution length} \(\linv(w)\). Their second conjecture reads, verbatim~\cite[Conj.~1.3]{HR}:

\begin{conjecture*}[Hart--Rowley, Conjecture~1.3]
Suppose Conjecture 1.2 holds, and let \(W\) be a finite Coxeter group. For each \(w \in W\), the involution length of \(w\) is at most the rank of \(W\).
\end{conjecture*}

The statement is the same, word for word, in the arXiv version and in the journal version of~\cite{HR}.

\subsection{What was known}
Hart and Rowley proved that Coxeter elements have the ancestor property~\cite[Prop.~2.2]{HR}, and they reported a computer check of both conjectures for the irreducible finite Coxeter groups of order below \(100{,}000\)~\cite[\S3]{HR}. Gaetz and Gao proved Conjecture~1.2 for every Coxeter group~\cite[Cor.~1.2]{GG}; so the hypothesis of Conjecture~1.3 holds. Levinstein then showed that Conjecture~1.3 is false~\cite{Lev}. His preprint, dated 2026-11-27, gives counterexamples in \(H_3\) (involution length \(4\), rank \(3\)), \(F_4\) (\(5>4\)) and \(E_6\) (\(8>6\)), and deduces a violation in \(E_7\) (involution length at least \(8\), rank \(7\)). The groups \(H_3\), \(F_4\) and \(E_6\) are within the range of the computer check reported in~\cite[\S3]{HR}; Levinstein notes that this check is incorrect for Conjecture~1.3 in these three groups~\cite{Lev}. A companion note in the same repository as this paper verifies Levinstein's \(H_3\) counterexample independently (file \file{../paper.tex}, in the parent directory of this paper).

So Conjecture~1.3 is false in general. This paper shows that it holds for the symmetric groups, that is, in type A.

\subsection{Results}
Let \(S_N\) be the group of permutations of \([N]=\{1,\dots,N\}\), with the adjacent transpositions \(s_i=(i\ \ i+1)\) as generators. It is the Coxeter group of type \(A_{N-1}\), of rank \(N-1\). The \emph{displacement} of \(w\in S_N\) is
\begin{equation}\label{eq:disp}
\disp(w)=\max_{1\le i\le N}\abs{w(i)-i}.
\end{equation}

\begin{theoremA}[Theorem~\ref{thm:A}]
For every \(N\ge1\) and every \(w\in S_N\),
\[
\linv(w)\le\disp(w)\le N-1=\rank(S_N).
\]
\end{theoremA}

So Conjecture~1.3 holds in type A (Corollary~\ref{cor:A}). Along the ancestor decomposition the displacement drops at every step (Corollary~\ref{cor:A}). Both inequalities are equalities for the Coxeter element \(s_1s_2\cdots s_{N-1}\) (Proposition~\ref{prop:attained}).

The group \(B_n\) of signed permutations sits inside the symmetric group on \(2n\) letters (Section~\ref{sec:typeB}).

\begin{theoremB}[Theorem~\ref{thm:B}]
For every \(n\ge1\) and every \(w\in B_n\), the involution length of \(w\) in \(B_n\) equals the involution length of \(w\) in \(S_{2n}\). Hence it is at most \(2n-1\).
\end{theoremB}

Theorem~B is \emph{not} Conjecture~1.3 for type B. The rank of \(B_n\) is \(n\), and \(2n-1>n\) for \(n\ge2\). Conjecture~1.3 for \(B_n\) remains open here. The analogous reduction is false for type D already in \(D_3\) (Remark~\ref{rem:D}).

\subsection{Idea of the proof}
Let \(\rem(w)=a(w)^{-1}w\) be what remains after the ancestor is removed. By induction on length, Theorem~A follows from one inequality: \(\disp(\rem(w))<\disp(w)\) for \(w\neq e\) (Proposition~\ref{prop:drop}). The remainder lies below \(w\) in the Bruhat order, and displacement cannot increase going down in that order; so \(\disp(\rem(w))\le\disp(w)\). To get a strict drop, take a position \(p\) where \(w\) has its largest displacement \(D\), say \(w(p)=p+D\). An explicit construction (Lemma~\ref{lem:K}) gives an involution prefix \(\iota\) of \(w\) that moves the value \(p+D\). The theorem of Gaetz and Gao puts \(\iota\) below the ancestor in Bruhat order. A complementation step then puts \(\rem(w)\) below \(\iota w\). Since \(\iota w\) does not have the value \(p+D\) at position \(p\), neither does \(\rem(w)\). Positions where \(w(p)=p-D\) are handled by the symmetry \(x\mapsto w_0xw_0\).

\subsection{What is proved how}
We keep three kinds of evidence apart.
\begin{itemize}
\item \emph{By hand.} Every theorem, proposition and lemma of this paper is proved by hand. This includes the theorem of Gaetz and Gao that we use (Theorem~\ref{thm:GG}); Appendix~\ref{app:GG} writes out their proof from textbook facts. The only inputs not proved here are the textbook facts listed in Section~\ref{sec:facts} and Fact~\ref{fact:B}, from Bj\"orner and Brenti~\cite{BB}.
\item \emph{By computer.} One exact implementation (two Python scripts) checks the main statements and the lemmas on all of \(S_N\) for \(N\le7\), and the type B and type D statements for \(n\le4\) (Section~\ref{sec:computer}). This is corroboration only. No step of any proof uses a computation.
\item \emph{In Lean.} Nothing in this paper is formalized (Section~\ref{sec:lean}).
\end{itemize}

\subsection{Credit}
The conjectures, the ancestor decomposition and the involution length are due to Hart and Rowley~\cite{HR}. The ancestor property of Coxeter elements and the involution length of Coxeter elements, which gives the value \(\linv(s_1s_2\cdots s_{N-1})=N-1\) in Proposition~\ref{prop:attained}, are due to Hart and Rowley~\cite[Prop.~2.2 and Prop.~2.5]{HR}. The uniqueness of the ancestor, the fact that the ancestor lies above every involution prefix in Bruhat order, and the complementation step (Lemma~\ref{lem:compl}) are due to Gaetz and Gao~\cite{GG}; Appendix~\ref{app:GG} is their argument, written out in our notation. The refutation of Conjecture~1.3 in general is due to Levinstein~\cite{Lev}. We make no claim of priority or novelty.

\section{Coxeter groups: definitions and the facts we use}\label{sec:coxeter}

\subsection{Definitions}
A \emph{Coxeter system} \((W,S)\) is a group \(W\) with a finite generating set \(S\) and the presentation
\begin{equation}\label{eq:presentation}
W=\bigl\langle S \bigm| (st)^{m(s,t)}=e \text{ for all } s,t\in S\bigr\rangle,
\end{equation}
where \(e\) is the identity, \(m(s,s)=1\), and \(m(s,t)=m(t,s)\in\{2,3,\dots\}\cup\{\infty\}\) for \(s\neq t\); there is no relation when \(m(s,t)=\infty\). The \emph{rank} is \(\abs{S}\). Each generator satisfies \(s^2=e\).

A \emph{word} is a finite sequence of generators; it represents the product of its letters. The \emph{length} \(\ell(w)\) is the least number of letters of a word representing \(w\). A word for \(w\) with \(\ell(w)\) letters is a \emph{reduced word} for \(w\). The set of \emph{reflections} is \(T=\{wsw^{-1}\mid w\in W,\ s\in S\}\). A generator \(s\) is a \emph{right descent} of \(w\) if \(\ell(ws)<\ell(w)\), and a \emph{left descent} if \(\ell(sw)<\ell(w)\).

\begin{definition}[weak orders]\label{def:weak}
For \(u,w\in W\):
\begin{itemize}
\item \(u\leR w\) (\emph{\(u\) is a prefix of \(w\)}) if \(\ell(w)=\ell(u)+\ell(u^{-1}w)\);
\item \(u\leL w\) (\emph{\(u\) is a suffix of \(w\)}) if \(\ell(w)=\ell(wu^{-1})+\ell(u)\).
\end{itemize}
\end{definition}

These are the right and left weak orders of~\cite[Ch.~3]{BB}. Hart and Rowley define a prefix through reduced words: \(u\) is a prefix of \(w\) if a reduced word of \(w\) begins with a reduced word of \(u\)~\cite{HR}. The two definitions agree. Indeed, if \(\ell(w)=\ell(u)+\ell(u^{-1}w)\), a reduced word for \(u\) followed by one for \(u^{-1}w\) is a word for \(w\) with \(\ell(w)\) letters, hence reduced. Conversely, if a reduced word for \(w\) begins with a reduced word for \(u\), the remaining letters represent \(u^{-1}w\), so \(\ell(u^{-1}w)\le\ell(w)-\ell(u)\); the reverse inequality follows from \(w=u\cdot u^{-1}w\). Hart and Rowley call this order the weak (left) Bruhat order; only the name differs. If \(s\) is the first letter of a reduced word for \(w\), then \(s\leR w\).

\begin{definition}[Bruhat order]\label{def:bruhat}
\(x\leB y\) if there is a chain \(x=x_0,x_1,\dots,x_m=y\) with \(m\ge0\), \(x_i=x_{i-1}t_i\) for some \(t_i\in T\), and \(\ell(x_i)>\ell(x_{i-1})\).
\end{definition}

This is the Bruhat order of~\cite[Ch.~2]{BB}. Since \(T\) is closed under conjugation, \(x_{i-1}t_i=(x_{i-1}t_ix_{i-1}^{-1})x_{i-1}\), so it does not matter whether the reflections multiply on the right or on the left. Along a chain the length strictly increases. Hence
\begin{equation}\label{eq:I2}
x\leB y \text{ and } x\neq y \implies \ell(x)<\ell(y).
\end{equation}
In particular \(\leB\) is antisymmetric, so it is a partial order.

An \emph{involution} is an element \(u\) with \(u^2=e\) and \(u\neq e\). The identity is not an involution.

\subsection{Textbook facts}\label{sec:facts}
We use the following standard facts and nothing else from outside this paper.

\begin{fact}[basic properties of length; {\cite[Ch.~1]{BB}}]\label{fact:length}
For \(x,y\in W\) and \(s\in S\): \(\ell(x^{-1})=\ell(x)\); \(\ell(xy)\le\ell(x)+\ell(y)\); and \(\ell(xs)=\ell(x)\pm1\), \(\ell(sx)=\ell(x)\pm1\).
\end{fact}

\begin{fact}[lifting property; {\cite[Ch.~2]{BB}}]\label{fact:lifting}
Let \(u\leB w\) with \(u\neq w\), and let \(s\in S\) be a right descent of \(w\) but not of \(u\). Then \(u\leB ws\) and \(us\leB w\).
\end{fact}

\begin{fact}[subword property; {\cite[Ch.~2]{BB}}]\label{fact:subword}
Let \(s_1s_2\cdots s_m\) be a reduced word for \(y\). Then \(u\leB y\) if and only if some subword \(s_{i_1}s_{i_2}\cdots s_{i_q}\), with \(i_1<i_2<\dots<i_q\), is a reduced word for \(u\).
\end{fact}

\begin{fact}[the symmetric group; {\cite[Ch.~1]{BB}}]\label{fact:SN}
\((S_N,\{s_1,\dots,s_{N-1}\})\) is a Coxeter system of type \(A_{N-1}\), and \(\ell(w)\) is the number of inversions of \(w\), that is, of pairs \(i<j\) with \(w(i)>w(j)\).
\end{fact}

The reflections of \(S_N\) are the conjugates of the adjacent transpositions, that is, all transpositions \((i\ j)\). The facts about the groups of types B and D are recalled in Section~\ref{sec:typeB}. Some books state the lifting property with left descents instead of right ones. The two versions are equivalent by Lemma~\ref{lem:inverse}(2) below, so either can be taken as the input.

\subsection{Elementary consequences}

\begin{lemma}[inversion]\label{lem:inverse}
For \(g,h,u,w\in W\):
\begin{enumerate}
\item \(u\leR w\) if and only if \(u^{-1}\leL w^{-1}\);
\item \(g\leB h\) if and only if \(g^{-1}\leB h^{-1}\).
\end{enumerate}
\end{lemma}

\begin{proof}
(1) \(u^{-1}\leL w^{-1}\) means \(\ell(w^{-1})=\ell(w^{-1}u)+\ell(u^{-1})\). By Fact~\ref{fact:length} this is the equation \(\ell(w)=\ell(u^{-1}w)+\ell(u)\), which defines \(u\leR w\).
(2) A chain step \(x\to xt\) with \(t\in T\) and \(\ell(xt)>\ell(x)\) becomes \(x^{-1}\to tx^{-1}=x^{-1}(xtx^{-1})\), with \(xtx^{-1}\in T\) and the same lengths. So inversion maps chains to chains. It is its own inverse, which gives both directions.
\end{proof}

We use the lifting property in three forms.

\begin{lemma}[forms of the lifting property]\label{lem:lift}
Let \(u\leB w\) and let \(s\in S\).
\begin{enumerate}
\item[(R1)] If \(\ell(ws)<\ell(w)\) and \(\ell(us)<\ell(u)\), then \(us\leB ws\).
\item[(R2)] If \(\ell(ws)<\ell(w)\) and \(\ell(us)>\ell(u)\), then \(u\leB ws\).
\item[(L1)] If \(\ell(sw)<\ell(w)\) and \(\ell(su)<\ell(u)\), then \(su\leB sw\).
\end{enumerate}
\end{lemma}

\begin{proof}
(R2) Here \(u\neq w\), since \(s\) is a right descent of \(w\) but not of \(u\). Fact~\ref{fact:lifting} gives \(u\leB ws\).

(R1) If \(u=w\) there is nothing to prove. Otherwise, first, \(us\leB u\): indeed \(u=(us)s\) with \(s\in T\) and \(\ell(u)>\ell(us)\), a chain step. So \(us\leB w\). Next, \(us\neq w\), because \(\ell(us)<\ell(u)\le\ell(w)\) by~\eqref{eq:I2}. Finally \(s\) is a right descent of \(w\) and not of \(us\), since \(\ell(us\cdot s)=\ell(u)>\ell(us)\). Fact~\ref{fact:lifting} for the pair \(us,w\) gives \(us\leB ws\).

(L1) Apply (R1) to \(u^{-1}\leB w^{-1}\) (Lemma~\ref{lem:inverse}(2)): \(s\) is a right descent of \(w^{-1}\) and of \(u^{-1}\), because \(\ell(w^{-1}s)=\ell(sw)\) and \(\ell(u^{-1}s)=\ell(su)\) by Fact~\ref{fact:length}. So \(u^{-1}s\leB w^{-1}s\), and inverting again gives \(su\leB sw\).
\end{proof}

\begin{lemma}[suffixes and prefixes lie below]\label{lem:suffix}
Let \(W\) be any Coxeter group and \(r,w\in W\).
\begin{enumerate}
\item If \(r\leL w\), then \(r\leB w\).
\item If \(r\leR w\), then \(r\leB w\).
\end{enumerate}
\end{lemma}

\begin{proof}
(1) Write \(w=cr\) with \(\ell(w)=\ell(c)+\ell(r)\). We induct on \(\ell(c)\). If \(c=e\), then \(r=w\). Otherwise write \(c=sc'\) with \(s\in S\) and \(\ell(c')=\ell(c)-1\); here \(s\) is the first letter of a reduced word for \(c\). By Fact~\ref{fact:length}, \(\ell(c'r)\le\ell(c')+\ell(r)=\ell(w)-1\) and \(\ell(w)=\ell(s\cdot c'r)\le\ell(c'r)+1\). So \(\ell(c'r)=\ell(w)-1=\ell(c')+\ell(r)\), and the induction hypothesis gives \(r\leB c'r\). Finally \(w=s(c'r)=(c'r)t\) with \(t=(c'r)^{-1}s(c'r)\in T\) and \(\ell(w)>\ell(c'r)\). So \(c'r\leB w\), and \(r\leB w\).

(2) By Lemma~\ref{lem:inverse}(1), \(r^{-1}\leL w^{-1}\). By (1), \(r^{-1}\leB w^{-1}\). By Lemma~\ref{lem:inverse}(2), \(r\leB w\).
\end{proof}

The next lemma is the first step of Gaetz and Gao's proof of their Lemma~3.1~\cite{GG}. We write it out because it is used twice, in Section~\ref{sec:typeA} and in Appendix~\ref{app:GG}.

\begin{lemma}[complementation; Gaetz--Gao]\label{lem:compl}
Let \((W,S)\) be a Coxeter system and \(u,v,w\in W\) with \(u\leR w\), \(v\leR w\) and \(u\leB v\). Then \(v^{-1}w\leB u^{-1}w\).
\end{lemma}

\begin{proof}
We induct on \(\ell(v)\). If \(\ell(v)=0\), then \(v=e\); \(u\leB e\) forces \(u=e\) by~\eqref{eq:I2}, and both sides equal \(w\). Let \(\ell(v)\ge1\) and choose \(s\in S\) with \(\ell(vs)=\ell(v)-1\). Put \(a=v^{-1}w\) and \(b=u^{-1}w\). Then \(\ell(w)=\ell(v)+\ell(a)=\ell(u)+\ell(b)\).

\emph{Claim.} If \(z\leR w\) and \(\ell(zs)=\ell(z)-1\), then \(zs\leR w\) and \(\ell(sz^{-1}w)=\ell(z^{-1}w)+1\). Indeed \(w=(zs)(sz^{-1}w)\), so by Fact~\ref{fact:length}
\[
\ell(w)\le\ell(zs)+\ell(sz^{-1}w)\le(\ell(z)-1)+(\ell(z^{-1}w)+1)=\ell(w).
\]
Both inequalities are equalities. Since \((zs)^{-1}w=sz^{-1}w\), this is the claim.

The claim for \(z=v\) gives \(vs\leR w\), \((vs)^{-1}w=sa\) and \(\ell(sa)=\ell(a)+1\).

\emph{Case 1: \(\ell(us)<\ell(u)\).} The claim for \(z=u\) gives \(us\leR w\), \((us)^{-1}w=sb\) and \(\ell(sb)=\ell(b)+1\). By (R1) of Lemma~\ref{lem:lift}, \(us\leB vs\). The induction hypothesis for the triple \((us,vs,w)\), with \(\ell(vs)<\ell(v)\), gives \(sa\leB sb\). Now \(s\) is a left descent of both \(sa\) and \(sb\). So (L1) gives \(s(sa)\leB s(sb)\), that is, \(a\leB b\).

\emph{Case 2: \(\ell(us)>\ell(u)\).} By (R2), \(u\leB vs\). The induction hypothesis for the triple \((u,vs,w)\) gives \(sa\leB b\). Also \(a\leB sa\), because \(sa=a\cdot(a^{-1}sa)\) with \(a^{-1}sa\in T\) and \(\ell(sa)>\ell(a)\). Hence \(a\leB b\).
\end{proof}

\subsection{The theorem of Gaetz and Gao}
The following theorem is Corollary~1.2 of~\cite{GG}, restated in our notation. It holds in every Coxeter group.

\begin{theorem}[Gaetz--Gao]\label{thm:GG}
Let \((W,S)\) be a Coxeter system and \(w\in W\), \(w\neq e\).
\begin{enumerate}
\item Among the involutions \(u\) with \(u\leR w\) there is exactly one of largest length.
\item Call it \(a(w)\). Every involution \(\iota\) with \(\iota\leR w\) satisfies \(\iota\leB a(w)\).
\end{enumerate}
\end{theorem}

The preprint~\cite{GG} has not been refereed. We therefore do not rely on it as an authority. Appendix~\ref{app:GG} writes out Gaetz and Gao's proof in full, in our notation. That proof uses only Facts~\ref{fact:length}--\ref{fact:subword}, the bound~\eqref{eq:I2}, and Lemmas~\ref{lem:inverse}--\ref{lem:compl}. None of these uses Theorem~\ref{thm:GG}, so there is no circularity. From now on Theorem~\ref{thm:GG} is used as proved.

\subsection{Ancestor and involution length}

\begin{definition}[{\cite{HR}}]\label{def:ancestor}
Let \(w\in W\), \(w\neq e\). The set \(\IP(w)\) of \emph{involution prefixes} of \(w\) is the set of involutions \(u\) with \(u\leR w\). It is not empty: the first letter of a reduced word for \(w\) is in it. The \emph{ancestor} \(a(w)\) is the unique element of \(\IP(w)\) of largest length (Theorem~\ref{thm:GG}(1)). The \emph{remainder} is \(\rem(w)=a(w)^{-1}w=a(w)w\).
\end{definition}

Since \(a(w)\leR w\), we have \(\ell(\rem(w))=\ell(w)-\ell(a(w))<\ell(w)\). This allows a definition by induction on length.

\begin{definition}[{\cite{HR}}]\label{def:linv}
The \emph{involution length} is defined by
\begin{equation}\label{eq:linv}
\linv(e)=0,\qquad \linv(w)=1+\linv(\rem(w))\quad (w\neq e).
\end{equation}
\end{definition}

This is the number \(k\) of factors of the ancestor decomposition \(w=a_1a_2\cdots a_k\) of Hart and Rowley, where \(a_i=a(a_ia_{i+1}\cdots a_k)\). Indeed \(a_1=a(w)\), and \(a_2\cdots a_k=\rem(w)\) is decomposed in the same way. If \(w\) is an involution, then \(w\in\IP(w)\) is its own longest involution prefix, so \(a(w)=w\), \(\rem(w)=e\) and \(\linv(w)=1\).

In this notation, Conjecture~1.3 of~\cite{HR} asserts \(\linv(w)\le\abs{S}\) for every element \(w\) of every finite Coxeter system \((W,S)\).

\section{The symmetric group}\label{sec:SN}

\subsection{Conventions}
Permutations of \([N]\) are composed as functions: \((uv)(i)=u(v(i))\). We write \(w\in S_N\) in one-line notation \(w(1)w(2)\cdots w(N)\). ``The value \(b\) stands at position \(j\)'' means \(w(j)=b\). Left multiplication acts on values, right multiplication on positions. The one-line notation of \((i\ j)w\) is that of \(w\) with the \emph{values} \(i\) and \(j\) exchanged. The one-line notation of \(w(i\ j)\) is that of \(w\) with the entries at the \emph{positions} \(i\) and \(j\) exchanged.

For \(w\in S_N\) define two sets of pairs:
\[
\PInv(w)=\{(a,b)\mid a<b,\ w(a)>w(b)\},\qquad
\VInv(w)=\{(a,b)\mid a<b,\ w^{-1}(a)>w^{-1}(b)\}.
\]
\(\PInv(w)\) is the set of inversions of \(w\), recorded as pairs of positions. \(\VInv(w)\) is the set of pairs of values \(a<b\) such that \(b\) stands to the left of \(a\). By definition \(\VInv(w)=\PInv(w^{-1})\). The map \((i,j)\mapsto(w(j),w(i))\) is a bijection from \(\PInv(w)\) onto \(\VInv(w)\). So by Fact~\ref{fact:SN},
\begin{equation}\label{eq:lengthSN}
\abs{\PInv(w)}=\abs{\VInv(w)}=\ell(w).
\end{equation}

Further notation. A position \(p\) with \(w(p)-p=d>0\) is an \emph{excedance of size \(d\)}; a position with \(p-w(p)=d>0\) is a \emph{drop of size \(d\)}. The \emph{longest element} \(w_0\) is the permutation \(i\mapsto N+1-i\). Put \(\varphi(x)=w_0xw_0\), so that \(\varphi(x)(i)=N+1-x(N+1-i)\). The \emph{rank matrix} of \(x\in S_N\) has entries
\[
x[i,k]=\card\{j\le i\mid x(j)\ge k\}\qquad(i,k\in[N]).
\]
Everything in this section uses only the order of \([N]\). So it applies to the permutations of any finite totally ordered set; this is used in Section~\ref{sec:typeB}.

\subsection{Inversions and weak order}

\begin{lemma}[inversions of a product]\label{lem:product}
Let \(x,y\in S_N\) and let \(i<j\) be positions. Write \(\{y(i),y(j)\}=\{a,b\}\) with \(a<b\). Then \((i,j)\in\PInv(xy)\) if and only if \((a,b)\) lies in exactly one of \(\PInv(x)\) and \(\VInv(y)\). Consequently
\begin{equation}\label{eq:product}
\ell(xy)=\ell(x)+\ell(y)-2\abs{\PInv(x)\cap\VInv(y)}.
\end{equation}
\end{lemma}

\begin{proof}
The pair \((i,j)\) is in \(\PInv(xy)\) when \(x(y(i))>x(y(j))\). The positions of \(a\) and \(b\) in \(y\) are \(i\) and \(j\) in some order. The condition \((a,b)\in\VInv(y)\) means \(y(i)=b\) and \(y(j)=a\). If \((a,b)\notin\VInv(y)\), then \(y(i)=a\), \(y(j)=b\), and the condition reads \(x(a)>x(b)\), that is, \((a,b)\in\PInv(x)\). If \((a,b)\in\VInv(y)\), the condition reads \(x(b)>x(a)\), that is, \((a,b)\notin\PInv(x)\). This proves the first statement.

The map \(\{i,j\}\mapsto\{y(i),y(j)\}\) is a bijection on the two-element subsets of \([N]\). So \(\abs{\PInv(xy)}\) is the size of the symmetric difference of \(\PInv(x)\) and \(\VInv(y)\). By~\eqref{eq:lengthSN} this gives~\eqref{eq:product}.
\end{proof}

\begin{lemma}[prefix criterion]\label{lem:prefixcrit}
For \(u,w\in S_N\): \(u\leR w\) if and only if \(\VInv(u)\subseteq\VInv(w)\). In particular, an involution \(\iota\) satisfies \(\iota\leR w\) if and only if \(\PInv(\iota)\subseteq\VInv(w)\).
\end{lemma}

\begin{proof}
Apply Lemma~\ref{lem:product} to \(x=u^{-1}\) and \(y=w\). Since \(\PInv(u^{-1})=\VInv(u)\) and \(\ell(u^{-1})=\ell(u)\),
\[
\ell(u^{-1}w)=\ell(u)+\ell(w)-2\abs{\VInv(u)\cap\VInv(w)}.
\]
So \(\ell(u^{-1}w)=\ell(w)-\ell(u)\) holds exactly when \(\abs{\VInv(u)\cap\VInv(w)}=\ell(u)=\abs{\VInv(u)}\), that is, when \(\VInv(u)\subseteq\VInv(w)\). For an involution, \(\VInv(\iota)=\PInv(\iota^{-1})=\PInv(\iota)\).
\end{proof}

\begin{lemma}[transpositions and length]\label{lem:transp}
Let \(x\in S_N\) and \(i<j\). Then \(\ell(x(i\ j))>\ell(x)\) if and only if \(x(i)<x(j)\).
\end{lemma}

\begin{proof}
Let \(y=x(i\ j)\), so \(y(i)=x(j)\), \(y(j)=x(i)\), and \(y(c)=x(c)\) for \(c\notin\{i,j\}\). Assume \(x(i)<x(j)\). We compare the inversions of \(x\) and \(y\) pair by pair. Pairs of positions disjoint from \(\{i,j\}\) contribute equally. For \(c<i\), the two pairs \((c,i)\) and \((c,j)\) together contribute \([x(c)>x(i)]+[x(c)>x(j)]\) to both \(x\) and \(y\), where \([P]\) is \(1\) if \(P\) holds and \(0\) otherwise; the same holds for \(c>j\). For \(i<c<j\), the pairs \((i,c)\) and \((c,j)\) contribute \([x(i)>x(c)]+[x(c)>x(j)]\) to \(x\) and \([x(j)>x(c)]+[x(c)>x(i)]\) to \(y\). Since \(x(j)>x(i)\), the second sum is at least the first. The pair \((i,j)\) contributes \(0\) to \(x\) and \(1\) to \(y\). Hence \(\ell(y)\ge\ell(x)+1\).

If instead \(x(i)>x(j)\), then \(y(i)<y(j)\) and \(x=y(i\ j)\). The case just proved gives \(\ell(x)>\ell(y)\).
\end{proof}

\subsection{Bruhat order and the symmetry \texorpdfstring{\(\varphi\)}{phi}}

\begin{lemma}[Bruhat order and rank matrices]\label{lem:rankmatrix}
If \(x\leB y\) in \(S_N\), then \(x[i,k]\le y[i,k]\) for all \(i,k\in[N]\).
\end{lemma}

\begin{proof}
It suffices to treat one step \(y=x(i\ j)\) with \(i<j\) and \(\ell(y)>\ell(x)\). By Lemma~\ref{lem:transp}, \(x(i)<x(j)\). For \(m<i\) or \(m\ge j\) the sets \(\{x(1),\dots,x(m)\}\) and \(\{y(1),\dots,y(m)\}\) are equal, so \(y[m,k]=x[m,k]\). For \(i\le m<j\) the second set is the first with \(x(i)\) replaced by the larger value \(x(j)\). So \(y[m,k]-x[m,k]=[x(j)\ge k]-[x(i)\ge k]\ge0\).
\end{proof}

Only this implication is needed. (Its converse is also true, but we do not use it.)

\begin{lemma}[the symmetry \(\varphi\)]\label{lem:phi}
The map \(\varphi(x)=w_0xw_0\) is a group automorphism of \(S_N\) with \(\varphi(s_i)=s_{N-i}\). It preserves length, the set \(T\) of reflections, involutions, \(\leR\) and \(\leB\), and \(\disp(\varphi(x))=\disp(x)\). For \(w\neq e\), \(a(\varphi(w))=\varphi(a(w))\) and \(\rem(\varphi(w))=\varphi(\rem(w))\). For every position \(p\) and integer \(d\),
\begin{equation}\label{eq:phiflip}
x(p)=p-d\iff \varphi(x)(N+1-p)=(N+1-p)+d.
\end{equation}
\end{lemma}

\begin{proof}
Conjugation by \(w_0\) sends \((i\ \ i+1)\) to \((N+1-i\ \ N-i)=s_{N-i}\). So \(\varphi\) is an automorphism that permutes \(S\). It maps reduced words to reduced words, so it preserves length. It preserves \(T\), the set of conjugates of \(S\), and the order of every element. Since \(\leR\) and \(\leB\) are defined from the group law, \(\ell\) and \(T\), they are preserved. So \(\varphi\) maps \(\IP(w)\) onto \(\IP(\varphi(w))\), preserving length. By the uniqueness in Theorem~\ref{thm:GG}(1), \(a(\varphi(w))=\varphi(a(w))\). Then \(\rem(\varphi(w))=\varphi(a(w))^{-1}\varphi(w)=\varphi(\rem(w))\). From \(\varphi(x)(i)=N+1-x(N+1-i)\) we get \(\varphi(x)(N+1-p)-(N+1-p)=p-x(p)\). This gives~\eqref{eq:phiflip}, and, taking absolute values and the maximum over \(p\), \(\disp(\varphi(x))=\disp(x)\).
\end{proof}

\subsection{Inversions of an involution}
The last lemma of this section is used only in the hand computations of Sections~\ref{sec:typeB} and~\ref{sec:small}.

\begin{lemma}\label{lem:invinv}
Let \(\iota\in S_N\) be an involution.
\begin{enumerate}
\item If \((a\ b)\) with \(a<b\) is a \(2\)-cycle of \(\iota\), then \((a,b)\in\PInv(\iota)\).
\item If \((a\ b)\) with \(a<b\) is a \(2\)-cycle of \(\iota\) and \(a<c<b\), then: \(\iota(c)<b\) implies \((a,c)\in\PInv(\iota)\), and \(\iota(c)>a\) implies \((c,b)\in\PInv(\iota)\).
\item If every \(2\)-cycle of \(\iota\) has the form \((i\ \ i+1)\), then \(\PInv(\iota)\) is the set of the corresponding pairs \((i,i+1)\).
\end{enumerate}
\end{lemma}

\begin{proof}
(1) \(\iota(a)=b>a=\iota(b)\). (2) In the first case \(\iota(a)=b>\iota(c)\); in the second \(\iota(c)>a=\iota(b)\). (3) Here \(\iota(i)\le i+1\) and \(\iota(j)\ge j-1\) for all \(i,j\). If \(i<j\) and \(\iota(i)>\iota(j)\), then \(i+1\ge\iota(i)>\iota(j)\ge j-1\), so \(j=i+1\); and \(\iota(i)>\iota(i+1)\ge i\) forces \(\iota(i)=i+1\). Conversely each \(2\)-cycle gives an inversion by (1).
\end{proof}

\section{Type A}\label{sec:typeA}

\subsection{Displacement and Bruhat order}

\begin{lemma}\label{lem:upset}
Let \(d\ge1\). For \(1\le i\le N-d\) put
\[
C_i^d=\{x\in S_N\mid x[i,i+d]\ge1\}=\{x\in S_N\mid x(j)\ge i+d\text{ for some }j\le i\}.
\]
\begin{enumerate}
\item \(C_i^d\) is an up-set of the Bruhat order: if \(x\in C_i^d\) and \(x\leB y\), then \(y\in C_i^d\).
\item \(x\) has an excedance of size at least \(d\) if and only if \(x\in C_i^d\) for some \(i\).
\item If \(\disp(x)\le d\) and \(1\le i\le N-d\), then \(x\in C_i^d\) if and only if \(x(i)=i+d\).
\item Displacement is monotone: \(x\leB y\) implies \(\disp(x)\le\disp(y)\).
\end{enumerate}
\end{lemma}

\begin{proof}
(1) By Lemma~\ref{lem:rankmatrix}, \(y[i,i+d]\ge x[i,i+d]\ge1\).

(2) If \(x(j)-j\ge d\), then \(j\le N-d\) and \(x\in C_j^d\). Conversely, if \(j\le i\) and \(x(j)\ge i+d\), then \(x(j)-j\ge i+d-j\ge d\).

(3) If \(x(i)=i+d\), then \(x\in C_i^d\) (take \(j=i\)). Conversely, let \(j\le i\) with \(x(j)\ge i+d\). Then \(d\ge x(j)-j\ge i+d-j\ge d\). All are equalities, so \(j=i\) and \(x(i)=i+d\).

(4) Let \(E_d\) be the set of permutations with an excedance of size at least \(d\). By (2) and (1), \(E_d\) is a union of up-sets, hence an up-set. By~\eqref{eq:phiflip}, \(x\) has a drop of size at least \(d\) if and only if \(\varphi(x)\in E_d\). Since \(\varphi\) preserves Bruhat order (Lemma~\ref{lem:phi}), \(\varphi^{-1}(E_d)\) is an up-set. So \(\{x\mid\disp(x)\ge d\}=E_d\cup\varphi^{-1}(E_d)\) is an up-set. Now let \(x\leB y\). If \(d=\disp(x)\ge1\), then \(y\) lies in this up-set, so \(\disp(y)\ge d\). If \(\disp(x)=0\), the inequality is clear.
\end{proof}

\subsection{An explicit involution prefix}
The next lemma is the combinatorial heart of the proof. Given a largest excedance of \(w\), it builds an involution prefix of \(w\) that moves the excedance's value.

\begin{lemma}\label{lem:K}
Let \(w\in S_N\) with \(D=\disp(w)\ge1\), and let \(p\) be a position with \(w(p)=p+D\). Put \(v=w(p)\). Then there is an involution \(\iota\in S_N\) with \(\iota\leR w\) and \(\iota(v)<v\).

Explicitly, let \(L=w(\{1,\dots,p\})\) and \(R=\{1,\dots,v\}\setminus L\). For \(1\le c\le v\) put
\begin{equation}\label{eq:h}
h(c)=\card\bigl(L\cap[c,v]\bigr)-\card\bigl(R\cap[c,v]\bigr).
\end{equation}
Let \(c_0\) be the largest \(c\) with \(h(c)=0\), and \(J=[c_0,v]\). Let \(l_1<\dots<l_k\) be the elements of \(L\cap J\) and \(r_1<\dots<r_k\) those of \(R\cap J\). Then \(\iota\) is the product of the transpositions \((r_i\ l_i)\), \(1\le i\le k\).
\end{lemma}

\begin{proof}
We first show that \(c_0\) exists and that \(L\cap J\) and \(R\cap J\) have the same size, so that \(\iota\) is well defined. Then we show that \(\iota\) is an involution moving \(v\) down, and finally that every inversion of \(\iota\) is an inversion of \(w\).

\emph{(a) \(L\subseteq\{1,\dots,v\}\), \(\abs{L}=p\), \(\abs{R}=D\) and \(v\in L\).} Let \(b>v\) be a value and \(j\) its position. Then \(b-j\le D\), so \(j\ge b-D>v-D=p\). So the positions \(1,\dots,p\) hold only values at most \(v\), and \(L\subseteq\{1,\dots,v\}\). Since \(w\) is injective, \(\abs{L}=p\), hence \(\abs{R}=v-p=D\). Also \(v=w(p)\in L\). Thus \(\{1,\dots,v\}\) is the disjoint union of \(L\) and \(R\): an element of \(\{1,\dots,v\}\) lies in \(L\) if its position is at most \(p\), and in \(R\) if its position is greater than \(p\).

\emph{(b) Every \(r\in R\) satisfies \(r\ge p+1-D\).} The position \(j\) of \(r\) is at least \(p+1\), and \(j-r\le D\).

\emph{(c) \(c_0\) exists, \(c_0<v\), and \(h(c)\ge1\) for \(c_0<c\le v\).} We have \(h(v)=1\) because \(v\in L\). For \(2\le c\le v\), \(h(c-1)-h(c)=\pm1\), because \(c-1\) lies in exactly one of \(L\) and \(R\). Put \(c_1=\max(1,p+1-D)\). By (b), and since \(v\notin R\), we have \(R\subseteq[c_1,v-1]\).
\begin{itemize}
\item If \(p+1-D\ge1\), then \(c_1=p+1-D\) and \([c_1,v]\) has \(v-c_1+1=2D\) elements. Of these, \(D\) lie in \(R\), hence \(D\) lie in \(L\). So \(h(c_1)=0\).
\item If \(p+1-D<1\), then \(c_1=1\) and \(h(1)=\abs{L}-\abs{R}=p-D<0\). Since \(h(v)=1>0\) and \(h\) changes by \(\pm1\) at each step, \(h\) vanishes somewhere strictly between \(1\) and \(v\).
\end{itemize}
In both cases \(h\) has a zero in \([1,v]\); let \(c_0\) be the largest. Then \(c_0<v\) because \(h(v)=1\). For \(c_0<c\le v\) we have \(h(c)\neq0\). Since \(h(v)=1\) and \(h\) moves in steps of \(\pm1\), it cannot become negative on \((c_0,v]\) without vanishing there; so \(h(c)\ge1\).

\emph{(d) \(\abs{L\cap J}=\abs{R\cap J}=k\ge1\), \(l_k=v\), and \(r_i<l_i\) for all \(i\).} The two sizes are equal because \(h(c_0)=0\). Also \(k\ge1\) and \(l_k=v\), because \(v=\max J\) lies in \(L\). For \(c_0\le c\le v-1\), taking complements inside \(J\) gives
\[
\card\bigl(R\cap[c_0,c]\bigr)-\card\bigl(L\cap[c_0,c]\bigr)=\bigl(k-\card(R\cap[c+1,v])\bigr)-\bigl(k-\card(L\cap[c+1,v])\bigr)=h(c+1)\ge1.
\]
For \(i<k\), apply this with \(c=l_i\le v-1\). Since \(\card(L\cap[c_0,l_i])=i\), we get \(\card(R\cap[c_0,l_i])\ge i+1\). Hence \(r_i\le l_i\), and \(r_i\neq l_i\) since \(L\cap R=\emptyset\). For \(i=k\): \(r_k\in J\) and \(r_k\neq v\), so \(r_k<v=l_k\).

\emph{(e) \(\iota\) is an involution with \(\iota(v)<v\).} The \(2k\) numbers \(l_i,r_i\) are distinct. So \(\iota\) is a product of \(k\ge1\) disjoint transpositions, and has order \(2\). Also \(\iota(v)=\iota(l_k)=r_k<v\).

\emph{(f) Every \((a,b)\in\PInv(\iota)\) has \(a\in R\cap J\) and \(b\in L\cap J\).} The permutation \(\iota\) fixes every point outside \(J\) and maps \(J\) onto \(J\). Let \(a<b\) with \(\iota(a)>\iota(b)\).
\begin{itemize}
\item If \(a<c_0\), then \(\iota(a)=a\). If \(b\in J\), then \(\iota(b)\ge c_0>a\); if \(b\notin J\), then \(\iota(b)=b>a\). Either way there is no inversion.
\item If \(a>v\), then \(b>v\), both are fixed, and there is no inversion.
\item If \(a\in J\) and \(b\notin J\), then \(b>v\), so \(\iota(b)=b>v\ge\iota(a)\). No inversion.
\end{itemize}
So \(a,b\in J\). If \(a=l_i\) and \(b=l_j\), then \(i<j\) and \(\iota(a)=r_i<r_j=\iota(b)\). If \(a=r_i\) and \(b=r_j\), then \(i<j\) and \(\iota(a)=l_i<l_j=\iota(b)\). If \(a=l_i\) and \(b=r_j\), then \(r_i<l_i<r_j\) gives \(i<j\), so \(\iota(a)=r_i<l_i<l_j=\iota(b)\). None of these is an inversion. The only remaining case is \(a\in R\cap J\) and \(b\in L\cap J\).

\emph{(g) \(\iota\leR w\).} Let \((a,b)\in\PInv(\iota)\). By (f), \(a<b\), \(a\in R\) and \(b\in L\). By (a), the position of \(b\) in \(w\) is at most \(p\) and the position of \(a\) is greater than \(p\). So \(w^{-1}(a)>w^{-1}(b)\), that is, \((a,b)\in\VInv(w)\). Hence \(\PInv(\iota)\subseteq\VInv(w)\), and Lemma~\ref{lem:prefixcrit} gives \(\iota\leR w\).
\end{proof}

For example, for \(w=21\) in \(S_2\): \(D=1\), \(p=1\), \(v=2\), \(L=\{2\}\), \(R=\{1\}\), \(h(2)=1\), \(h(1)=0\), \(c_0=1\), and \(\iota=(1\ 2)=w\). For \(w=312\) in \(S_3\): \(p=1\), \(v=3\), \(L=\{3\}\), \(R=\{1,2\}\), \(h(3)=1\), \(h(2)=0\), \(c_0=2\), and \(\iota=(2\ 3)=132\).

Both hypotheses on \(p\) matter. That \(D\) is the \emph{largest} displacement is used in (a) and (b); both can fail for a smaller excedance. The lemma treats excedances only. Drops are handled by the symmetry \(\varphi\) in Step~3 below.

\subsection{The remainder has smaller displacement}

\begin{proposition}\label{prop:drop}
Let \(w\in S_N\), \(w\neq e\), \(D=\disp(w)\) and \(r=\rem(w)\). Then \(\disp(r)\le D-1\).
\end{proposition}

\begin{proof}
We show three things: \(\disp(r)\le D\); \(r\) has no excedance of size \(D\); and \(r\) has no drop of size \(D\). Note \(D\ge1\) since \(w\neq e\). Write \(a=a(w)\).

\emph{Step 1: \(\disp(r)\le D\).} Since \(a\leR w\), we have \(w=a\cdot r\) with \(\ell(w)=\ell(a)+\ell(r)\). So \(r\leL w\). By Lemma~\ref{lem:suffix}(1), \(r\leB w\). By Lemma~\ref{lem:upset}(4), \(\disp(r)\le\disp(w)\).

\emph{Step 2: \(r(p)\neq p+D\) for every position \(p\).} We prove this for an arbitrary \(w\in S_N\setminus\{e\}\), since Step~3 applies it to another element. Suppose \(r(p)=p+D\). Then \(p\le N-D\) and \(r\in C_p^D\). Since \(r\leB w\) (Step~1) and \(C_p^D\) is an up-set (Lemma~\ref{lem:upset}(1)), \(w\in C_p^D\). Since \(\disp(w)=D\), Lemma~\ref{lem:upset}(3) gives \(w(p)=p+D\). So Lemma~\ref{lem:K} applies to \(w\) and \(p\). It gives an involution \(\iota\leR w\) with \(\iota(p+D)\neq p+D\).
\begin{itemize}
\item By Theorem~\ref{thm:GG}(2), applied in \(S_N\) to \(w\neq e\) and the involution \(\iota\leR w\): \(\iota\leB a\).
\item By Lemma~\ref{lem:compl}, applied to the triple \((\iota,a,w)\) (we have \(\iota\leR w\), \(a\leR w\) and \(\iota\leB a\)): \(a^{-1}w\leB\iota^{-1}w\), that is, \(r\leB\iota w\).
\item Since \(C_p^D\) is an up-set and \(r\in C_p^D\), also \(\iota w\in C_p^D\).
\item Since \(\iota\leR w\) and \(\iota^{-1}=\iota\), we have \(w=\iota\cdot(\iota w)\) with \(\ell(w)=\ell(\iota)+\ell(\iota w)\). So \(\iota w\leL w\). By Lemma~\ref{lem:suffix}(1) and Lemma~\ref{lem:upset}(4), \(\disp(\iota w)\le D\).
\item By Lemma~\ref{lem:upset}(3), \((\iota w)(p)=p+D\). But \((\iota w)(p)=\iota(w(p))=\iota(p+D)\neq p+D\). This is a contradiction.
\end{itemize}

\emph{Step 3: \(r(p)\neq p-D\) for every position \(p\).} Suppose \(r(p)=p-D\). Let \(w'=\varphi(w)\). By Lemma~\ref{lem:phi}, \(w'\neq e\), \(\disp(w')=D\) and \(\rem(w')=\varphi(r)\). By~\eqref{eq:phiflip}, \(\varphi(r)(N+1-p)=(N+1-p)+D\). This contradicts Step~2 for the element \(w'\) at the position \(N+1-p\).

By Steps 1--3, \(\abs{r(p)-p}\le D-1\) for every \(p\).
\end{proof}

\subsection{The main theorem}

\begin{theorem}\label{thm:A}
For every \(N\ge1\) and every \(w\in S_N\): \(\linv(w)\le\disp(w)\le N-1\).
\end{theorem}

\begin{proof}
The second inequality holds because \(w(i)\) and \(i\) both lie in \([N]\), so \(\abs{w(i)-i}\le N-1\). The first we prove by induction on \(\ell(w)\). For \(w=e\) both sides are \(0\). For \(w\neq e\), let \(r=\rem(w)\). Then \(\ell(r)<\ell(w)\) (Section~\ref{sec:coxeter}). By~\eqref{eq:linv}, the induction hypothesis for \(r\), and Proposition~\ref{prop:drop},
\[
\linv(w)=1+\linv(r)\le1+\disp(r)\le1+(\disp(w)-1)=\disp(w).\qedhere
\]
\end{proof}

\begin{corollary}\label{cor:A}
Hart--Rowley Conjecture~1.3 holds for the Coxeter system \((S_N,\{s_1,\dots,s_{N-1}\})\) of type \(A_{N-1}\), for every \(N\ge1\): \(\linv(w)\le N-1=\rank(S_N)\) for all \(w\in S_N\). Moreover, if \(w\neq e\) has ancestor decomposition \(w=a_1\cdots a_k\) and \(w_i=a_i\cdots a_k\), then
\[
\disp(w_1)>\disp(w_2)>\dots>\disp(w_k)\ge1.
\]
\end{corollary}

\begin{proof}
The first statement is Theorem~\ref{thm:A}. For the second, \(w_{i+1}=\rem(w_i)\) for \(i<k\), so Proposition~\ref{prop:drop} applies to each \(w_i\) with \(i<k\). Finally \(w_k\neq e\), so \(\disp(w_k)\ge1\).
\end{proof}

The bound is attained. Coxeter elements, that is, products of all the generators in some order, are studied by Hart and Rowley~\cite{HR}; the following case is due to them~\cite[Prop.~2.5 and the remarks after it]{HR}. We include a direct proof.

\begin{proposition}\label{prop:attained}
Let \(N\ge2\) and \(c=s_1s_2\cdots s_{N-1}\). Then \(\linv(c)=\disp(c)=N-1\).
\end{proposition}

\begin{proof}
We compute the ancestor decomposition of \(c\) explicitly: at each step the only involution prefix is a single generator.

For \(1\le k\le N\) let \(c_k=s_ks_{k+1}\cdots s_{N-1}\), so \(c_N=e\) and \(c_1=c\). We claim that the one-line notation of \(c_k\) is \(1,\dots,k-1,k+1,k+2,\dots,N,k\). For \(k=N\) this is \(e\). If it holds for \(k+1\), then \(c_k=s_kc_{k+1}\) is obtained from \(1,\dots,k,k+2,\dots,N,k+1\) by exchanging the values \(k\) and \(k+1\), which gives the claim for \(k\). So the only value that stands to the right of a larger value in \(c_k\) is \(k\), and
\[
\VInv(c_k)=\{(k,b)\mid k<b\le N\}.
\]
Also \(c(N)=1\), so \(\disp(c)=N-1\).

Let \(k\le N-1\). First, \(s_k\leR c_k\) by Lemma~\ref{lem:prefixcrit}, since \(\PInv(s_k)=\{(k,k+1)\}\subseteq\VInv(c_k)\). Now let \(\iota\) be any involution with \(\iota\leR c_k\). By Lemma~\ref{lem:prefixcrit}, \(\PInv(\iota)\subseteq\{(k,b)\mid b>k\}\). So no pair \((i,j)\) with \(i<j\) and \(i\neq k\) is an inversion of \(\iota\).
\begin{itemize}
\item For \(i<k\), we have \(\iota(i)<\iota(j)\) for all \(j>i\). By induction on \(i\), \(\iota(i)=i\): if \(\iota\) fixes every \(i'<i\), then \(\iota\) permutes \(\{i,\dots,N\}\), and \(\iota(i)\) is the smallest value taken there, so \(\iota(i)=i\).
\item Hence \(\iota\) permutes \(\{k,\dots,N\}\), and \(\iota(k+1)<\iota(k+2)<\dots<\iota(N)\).
\item If \(\iota(k)=k\), then \(\iota\) permutes \(\{k+1,\dots,N\}\) increasingly, so \(\iota=e\), which is not an involution. So \(\iota(k)=m>k\) and \(\iota(m)=k\).
\item The values taken on \(\{k+1,\dots,N\}\) are \(\{k,\dots,N\}\setminus\{m\}\), in increasing order. So \(\iota(k+1)=k\). With \(\iota(m)=k\) this gives \(m=k+1\).
\item Then \(\iota\) permutes \(\{k+2,\dots,N\}\) increasingly, so it fixes each of these points. Hence \(\iota=s_k\).
\end{itemize}
So \(\IP(c_k)=\{s_k\}\), \(a(c_k)=s_k\) and \(\rem(c_k)=s_kc_k=c_{k+1}\). Therefore \(\linv(c_k)=1+\linv(c_{k+1})\) for \(k\le N-1\), and \(\linv(c)=\linv(c_1)=N-1\).
\end{proof}

\section{Type B, and why type D is different}\label{sec:typeB}

\subsection{Setup}
Let \(n\ge1\) and \(I_n=\{-n,\dots,-1,1,\dots,n\}\) with its natural order. Let \(S(I_n)\) be the group of permutations of \(I_n\). The order-preserving bijection \(\lambda\colon I_n\to[2n]\), given by \(\lambda(x)=n+x\) for \(x>0\) and \(\lambda(x)=n+1+x\) for \(x<0\), identifies \(S(I_n)\) with \(S_{2n}\). All notions of Section~\ref{sec:SN} are taken in \(S(I_n)\cong S_{2n}\); we mark them with the letter A: the length \(\ell_A\), the prefix order \(\leR^A\), the ancestor \(a_A\), the remainder \(\rem_A\), the involution length \(\linv^A\) and the displacement \(\disp_A\). The longest element of \(S(I_n)\) is \(w_0(x)=-x\), and \(\varphi(z)=w_0zw_0\).

The group of \emph{signed permutations} is
\[
B_n=\{w\in S(I_n)\mid w(-x)=-w(x)\text{ for all }x\}=\{w\in S(I_n)\mid \varphi(w)=w\}.
\]
An element of \(B_n\) is determined by its \emph{window} \((w(1),\dots,w(n))\). Put \(s_0=(-1\ \ 1)\) and \(s_i=(i\ \ i+1)(-i\ \ -i-1)\) for \(1\le i\le n-1\). For the window define
\begin{align*}
\invB(w)&=\card\{1\le i<j\le n\mid w(i)>w(j)\},\qquad
\negB(w)=\card\{i\in[n]\mid w(i)<0\},\\
\nspB(w)&=\card\{1\le i<j\le n\mid w(i)+w(j)<0\}.
\end{align*}

\begin{fact}[{\cite[Ch.~8]{BB}}]\label{fact:B}
\((B_n,\{s_0,\dots,s_{n-1}\})\) is a Coxeter system of type \(B_n\), of rank \(n\), and its length function is \(\ell_B=\invB+\negB+\nspB\).
\end{fact}

We write \(\leR^B\), \(a_B\), \(\rem_B\) and \(\linv^B\) for the notions of Section~\ref{sec:coxeter} in this Coxeter system. A pair \((a,b)\) of elements of \(I_n\) with \(a<b\) is \emph{symmetric} if \(a=-b\). For a set \(P\) of pairs, \(F(P)\) denotes its subset of symmetric pairs.

Two facts about \(\negB\), for \(z\in B_n\):
\begin{itemize}
\item[(N1)] \(\negB(z)=\abs{F(\PInv(z))}=\abs{F(\VInv(z))}\). Indeed \((-c,c)\in\PInv(z)\) if and only if \(z(-c)>z(c)\), if and only if \(z(c)<0\); this gives the first equality. Applied to \(z^{-1}\in B_n\), it gives \(\abs{F(\VInv(z))}=\negB(z^{-1})\). Finally \(\negB(z^{-1})=\card\{d<0\mid z(d)>0\}=\card\{d>0\mid z(d)<0\}=\negB(z)\), by the substitution \(d=z^{-1}(c)\) and then \(z(-d)=-z(d)\).
\item[(N2)] In the labelling by \([2n]\), \(\negB(z)\) is the rank-matrix entry \(z[n,n+1]\). This entry counts the negative positions \(x\) with \(z(x)>0\), and \(\card\{c>0\mid z(-c)>0\}=\card\{c>0\mid z(c)<0\}\).
\end{itemize}

\subsection{Lengths and weak order}

\begin{lemma}\label{lem:lAlB}
For \(w\in B_n\): \(\ell_A(w)=2\ell_B(w)-\negB(w)\).
\end{lemma}

\begin{proof}
Count the inversions of \(w\) as a permutation of \(I_n\) by the signs of the two positions. Two positive positions \(i<j\) form an inversion when \(w(i)>w(j)\); there are \(\invB(w)\) of them. Two negative positions \(-j<-i\), with \(0<i<j\), form an inversion when \(w(-j)>w(-i)\), that is, when \(w(i)>w(j)\); again \(\invB(w)\) of them. A negative position \(-i\) and a positive position \(j\) form an inversion when \(w(-i)>w(j)\), that is, when \(w(i)+w(j)<0\). For \(i=j\) this gives \(\negB(w)\) pairs. For \(i\neq j\) the condition is symmetric in \(i\) and \(j\), so the ordered pairs give \(2\nspB(w)\). The total is \(2\invB+\negB+2\nspB=2\ell_B-\negB\), by Fact~\ref{fact:B}.
\end{proof}

\begin{lemma}\label{lem:Bweak}
For \(x,y\in B_n\) let \(P=\PInv(x)\cap\VInv(y)\). Then
\begin{equation}\label{eq:lBproduct}
\ell_B(xy)=\ell_B(x)+\ell_B(y)-\abs{P}-\abs{F(P)}.
\end{equation}
Consequently \(\ell_B(xy)=\ell_B(x)+\ell_B(y)\) if and only if \(P=\emptyset\), if and only if \(\ell_A(xy)=\ell_A(x)+\ell_A(y)\). For \(u,w\in B_n\): \(u\leR^B w\) if and only if \(u\leR^A w\). So the right weak order of \(B_n\) is the restriction of the right weak order of \(S_{2n}\).
\end{lemma}

\begin{proof}
By Lemma~\ref{lem:product}, \(\ell_A(xy)=\ell_A(x)+\ell_A(y)-2\abs{P}\). For \(\negB(xy)\) we use the first statement of Lemma~\ref{lem:product}. Since \(y\in B_n\), the bijection \(\{i,j\}\mapsto\{y(i),y(j)\}\) maps symmetric pairs to symmetric pairs, because \(\{y(-c),y(c)\}=\{-y(c),y(c)\}\); so does its inverse, given by \(y^{-1}\in B_n\). Therefore \(\abs{F(\PInv(xy))}\) is the size of the symmetric difference of \(F(\PInv(x))\) and \(F(\VInv(y))\). By (N1),
\[
\negB(xy)=\negB(x)+\negB(y)-2\abs{F(P)}.
\]
By Lemma~\ref{lem:lAlB}, \(2\ell_B(xy)=\ell_A(xy)+\negB(xy)=2\ell_B(x)+2\ell_B(y)-2\abs{P}-2\abs{F(P)}\), which is~\eqref{eq:lBproduct}. Since \(F(P)\subseteq P\), \(\abs{P}+\abs{F(P)}=0\) if and only if \(P=\emptyset\). By~\eqref{eq:product}, \(P=\emptyset\) if and only if \(\ell_A(xy)=\ell_A(x)+\ell_A(y)\). For the weak orders take \(x=u\) and \(y=u^{-1}w\).
\end{proof}

\subsection{Ancestors}

\begin{proposition}\label{prop:Bancestor}
Let \(w\in B_n\), \(w\neq e\). Then \(a_A(w)\in B_n\) and \(a_B(w)=a_A(w)\).
\end{proposition}

\begin{proof}
We show that the ancestor in \(S_{2n}\) lies in \(B_n\), is an involution prefix there, and is longer in \(\ell_B\) than every other involution prefix in \(B_n\).

Let \(a=a_A(w)\), the unique \(\ell_A\)-longest involution with \(a\leR^A w\) (Theorem~\ref{thm:GG}(1) for \(S_{2n}\)). By Lemma~\ref{lem:phi}, \(\varphi(a)\) is an involution with \(\varphi(a)\leR^A\varphi(w)=w\) and \(\ell_A(\varphi(a))=\ell_A(a)\). By uniqueness, \(\varphi(a)=a\), so \(a\in B_n\). By Lemma~\ref{lem:Bweak}, \(a\leR^B w\).

Let \(\iota\in B_n\) be any involution with \(\iota\leR^B w\). By Lemma~\ref{lem:Bweak}, \(\iota\leR^A w\). By Theorem~\ref{thm:GG}(2) for \(S_{2n}\), \(\iota\leB a\) in the Bruhat order of \(S_{2n}\). Hence \(\ell_A(\iota)\le\ell_A(a)\), with equality only if \(\iota=a\), by~\eqref{eq:I2}. Also \(\negB(\iota)=\iota[n,n+1]\le a[n,n+1]=\negB(a)\), by (N2) and Lemma~\ref{lem:rankmatrix}. By Lemma~\ref{lem:lAlB},
\[
2\ell_B(\iota)=\ell_A(\iota)+\negB(\iota)\le\ell_A(a)+\negB(a)=2\ell_B(a),
\]
with equality only if \(\ell_A(\iota)=\ell_A(a)\), that is, only if \(\iota=a\). So \(a\) is the unique \(\ell_B\)-longest involution prefix of \(w\) in \(B_n\).
\end{proof}

\begin{theorem}\label{thm:B}
For every \(n\ge1\) and every \(w\in B_n\): \(\linv^B(w)=\linv^A(w)\le\disp_A(w)\le2n-1\). In terms of the window,
\begin{equation}\label{eq:dispB}
\disp_A(w)=\max_{1\le i\le n}\delta_i,\qquad
\delta_i=\begin{cases}\abs{w(i)-i}&\text{if }w(i)>0,\\ i+\abs{w(i)}-1&\text{if }w(i)<0.\end{cases}
\end{equation}
\end{theorem}

\begin{proof}
We prove the equality by induction on \(\ell_B(w)\). For \(w=e\) both sides are \(0\). For \(w\neq e\), Proposition~\ref{prop:Bancestor} gives \(a_B(w)=a_A(w)=:a\in B_n\). So \(\rem_B(w)=a^{-1}w=\rem_A(w)\in B_n\), and \(\ell_B(\rem_B(w))<\ell_B(w)\). By the induction hypothesis,
\[
\linv^B(w)=1+\linv^B(a^{-1}w)=1+\linv^A(a^{-1}w)=\linv^A(w).
\]
The inequalities are Theorem~\ref{thm:A} with \(N=2n\). For~\eqref{eq:dispB}: the displacement at the position \(\lambda(i)\) is \(\abs{\lambda(w(i))-\lambda(i)}\). This equals \(\abs{w(i)-i}\) if \(i,w(i)>0\), and \(\abs{n+1+w(i)-n-i}=i+\abs{w(i)}-1\) if \(i>0>w(i)\). The positions \(\lambda(-i)\) give the same numbers, because \(w(-i)=-w(i)\) and \(\lambda(-x)=2n+1-\lambda(x)\).
\end{proof}

\begin{remark}[Theorem~\ref{thm:B} is not Conjecture~1.3 for type B]\label{rem:B}
The rank of \(B_n\) is \(n\), and Conjecture~1.3 for \(B_n\) asks for \(\linv^B\le n\). Theorem~\ref{thm:B} gives only \(2n-1\). Theorem~\ref{thm:A} cannot give more on its own: \(\disp_A\) takes the value \(2n-1\) on \(B_n\), for instance at \(w=-\mathrm{id}\), where \(\linv^B(w)=1\). The exact computation of Section~\ref{sec:computer} finds \(\max\linv^B=n\) for \(n\le4\); this is not used anywhere. We prove nothing about Conjecture~1.3 for \(B_n\).
\end{remark}

\subsection{Type D: the reduction fails}

\begin{remark}\label{rem:D}
Let \(D_n=\{w\in B_n\mid\negB(w)\text{ even}\}\), with generators \(s_0^D\), whose window is \((-2,-1,3,\dots,n)\), and \(s_1,\dots,s_{n-1}\). It is a Coxeter system of type \(D_n\), with length \(\ell_D=\invB+\nspB=\ell_B-\negB\)~\cite[Ch.~8]{BB}. The computation in the proof of Lemma~\ref{lem:Bweak} then gives
\[
\ell_D(xy)=\ell_D(x)+\ell_D(y)-\bigl(\abs{P}-\abs{F(P)}\bigr).
\]
So a product in \(D_n\) is \(\ell_D\)-additive exactly when \(P\) consists of symmetric pairs. This is weaker than \(P=\emptyset\). Hence \(u\leR^A w\) implies \(u\leR^D w\), but not conversely. Moreover \(D_n\) is not the fixed-point set of \(\varphi\). The following example shows that the ancestor in \(S_{2n}\) of an element of \(D_n\) need not lie in \(D_n\), and that the involution lengths can differ. So neither half of the argument for type B is available.

\emph{Example in \(D_3\).} Let \(w=s_2s_0^D\), with window \((-3,-1,2)\).
\begin{itemize}
\item \emph{In \(D_3\).} \(\ell_D(w)=\invB(w)+\nspB(w)=0+2=2\). So every reduced word for \(w\) has two letters, and its first letter \(s\) satisfies \(\ell_D(sw)=1\), that is, \(sw\) is a generator. Now \(s_2w=s_0^D\), while \(s_1w\) has window \((-3,-2,1)\) and \(s_0^Dw\) has window \((-3,2,-1)\); neither is a generator. So \(s_2s_0^D\) is the only reduced word for \(w\), and the prefixes of \(w\) are \(e\), \(s_2\) and \(w\). The element \(w\) is not an involution, since \(w(w(1))=w(-3)=-2\neq1\). Hence \(a_D(w)=s_2\), \(\rem_D(w)=s_0^D\), an involution, and \(\linv^D(w)=2\).
\item \emph{In \(S_6\).} In the labels \(\lambda\), \(w=246135\), and
\[
\VInv(w)=V=\{(1,2),(1,4),(1,6),(3,4),(3,6),(5,6)\}.
\]
Let \(\iota\) be an involution with \(\PInv(\iota)\subseteq V\). By Lemma~\ref{lem:invinv}(1), each \(2\)-cycle of \(\iota\) is a pair in \(V\). A \(2\)-cycle \((1\ 4)\) or \((1\ 6)\) is excluded by Lemma~\ref{lem:invinv}(2) with \(c=2\): \(\iota(2)\neq1\), so \(\iota(2)>1\), and \((2,4)\), respectively \((2,6)\), would be an inversion, but it is not in \(V\). A \(2\)-cycle \((3\ 6)\) is excluded with \(c=4\): if \(\iota(4)>3\), then \((4,6)\) is an inversion, not in \(V\); so \(\iota(4)<3\), and the \(2\)-cycle \((\iota(4)\ 4)\) is \((1\ 4)\), excluded, or \((2\ 4)\), not in \(V\). So the \(2\)-cycles of \(\iota\) are among \((1\ 2)\), \((3\ 4)\), \((5\ 6)\). By Lemma~\ref{lem:invinv}(3), every product of these is an involution prefix. Hence \(a_A(w)=214365\). Its window is \((-1,3,2)\), with one negative entry. So \(a_A(w)\in B_3\setminus D_3\).
\item \emph{Continuing in \(S_6\).} \(\rem_A(w)=135246\), with \(\VInv=\{(2,3),(2,5),(4,5)\}\). A \(2\)-cycle \((2\ 5)\) is excluded with \(c=3\): if \(\iota(3)>2\), then \((3,5)\) is an inversion, not in the set; so \(\iota(3)=1\), but \((1,3)\) is not in the set. So the ancestor is \(132546\), and the remainder is \(124356=s_3\in S_6\), an involution. Hence \(\linv^A(w)=3\neq2=\linv^D(w)\).
\end{itemize}
The exact computation of Section~\ref{sec:computer} finds this element as the shortest mismatch in \(D_3\). It also counts the mismatches: in \(D_3\), \(6\) elements have \(a_D\neq a_A\) and \(4\) have \(\linv^D\neq\linv^A\); in \(D_4\), \(55\) and \(33\). We claim nothing about Conjecture~1.3 for type D.
\end{remark}

\section{Small cases by hand}\label{sec:small}

These checks are not part of any proof. They illustrate the definitions and the conventions, and they can be compared with the computations of Section~\ref{sec:computer}.

\subsection{Procedure}
By Lemma~\ref{lem:prefixcrit}, the involution prefixes of \(w\in S_N\) are the involutions \(\iota\) with \(\VInv(\iota)\subseteq\VInv(w)\). So for each \(w\) we list \(\VInv(w)\), the pairs \(ab\) with \(a<b\) and \(b\) to the left of \(a\). We read off the involution prefixes from the list of involutions, take the longest one as \(a(w)\), and get \(\rem(w)=a(w)w\) by applying \(a(w)\) to the values of \(w\).

\subsection{\texorpdfstring{\(S_1\), \(S_2\) and \(S_3\)}{S1, S2 and S3}}
In \(S_1=\{e\}\), \(\linv=\disp=N-1=0\). In \(S_2\), \(\linv(s_1)=\disp(s_1)=N-1=1\). In \(S_3\) the involutions are \(213\) (with \(\VInv=\{12\}\)), \(132\) (\(\{23\}\)) and \(321\) (\(\{12,13,23\}\)). Table~\ref{tab:S3} lists all of \(S_3\).

\begin{table}[ht]
\centering
\small
\caption{The ancestor decomposition in \(S_3\).}\label{tab:S3}
\begin{tabular}{llllllll}
\toprule
\(w\) & \(\VInv(w)\) & involution prefixes & \(a(w)\) & \(\rem(w)\) & \(\linv\) & \(\disp(w)\) & \(\disp(\rem)\)\\
\midrule
123 & -- & -- & -- & -- & 0 & 0 & --\\
132 & 23 & 132 & 132 & 123 & 1 & 1 & 0\\
213 & 12 & 213 & 213 & 123 & 1 & 1 & 0\\
231 & 12 13 & 213 & 213 & 132 & 2 & 2 & 1\\
312 & 13 23 & 132 & 132 & 213 & 2 & 2 & 1\\
321 & 12 13 23 & 132 213 321 & 321 & 123 & 1 & 2 & 0\\
\bottomrule
\end{tabular}
\end{table}

The values \(0,1,2\) of \(\linv\) occur \(1\), \(3\) and \(2\) times. The maximum is \(2=N-1\). In every row \(\linv(w)\le\disp(w)\) and \(\disp(\rem(w))<\disp(w)\).

\subsection{\texorpdfstring{\(S_4\)}{S4}}
The nine involutions of \(S_4\), with their sets \(\VInv\) and lengths, are: \(2134\) (\(\{12\}\), length \(1\)); \(1324\) (\(\{23\}\), \(1\)); \(1243\) (\(\{34\}\), \(1\)); \(2143\) (\(\{12,34\}\), \(2\)); \(3214\) (\(\{12,13,23\}\), \(3\)); \(1432\) (\(\{23,24,34\}\), \(3\)); \(3412\) (\(\{13,14,23,24\}\), \(4\)); \(4231\) (\(\{12,13,14,24,34\}\), \(5\)); \(4321\) (all six pairs, \(6\)). Table~\ref{tab:S4} lists all of \(S_4\).

\begin{table}[ht]
\centering
\small
\caption{The ancestor decomposition in \(S_4\).}\label{tab:S4}
\begin{tabular}{llllllll}
\toprule
\(w\) & \(\VInv(w)\) & involution prefixes & \(a(w)\) & \(\rem(w)\) & \(\linv\) & \(\disp(w)\) & \(\disp(\rem)\)\\
\midrule
1234 & -- & -- & -- & -- & 0 & 0 & --\\
1243 & 34 & 1243 & 1243 & 1234 & 1 & 1 & 0\\
1324 & 23 & 1324 & 1324 & 1234 & 1 & 1 & 0\\
1342 & 23 24 & 1324 & 1324 & 1243 & 2 & 2 & 1\\
1423 & 24 34 & 1243 & 1243 & 1324 & 2 & 2 & 1\\
1432 & 23 24 34 & 1243 1324 1432 & 1432 & 1234 & 1 & 2 & 0\\
2134 & 12 & 2134 & 2134 & 1234 & 1 & 1 & 0\\
2143 & 12 34 & 1243 2134 2143 & 2143 & 1234 & 1 & 1 & 0\\
2314 & 12 13 & 2134 & 2134 & 1324 & 2 & 2 & 1\\
2341 & 12 13 14 & 2134 & 2134 & 1342 & 3 & 3 & 2\\
2413 & 12 14 34 & 1243 2134 2143 & 2143 & 1324 & 2 & 2 & 1\\
2431 & 12 13 14 34 & 1243 2134 2143 & 2143 & 1342 & 3 & 3 & 2\\
3124 & 13 23 & 1324 & 1324 & 2134 & 2 & 2 & 1\\
3142 & 13 23 24 & 1324 & 1324 & 2143 & 2 & 2 & 1\\
3214 & 12 13 23 & 1324 2134 3214 & 3214 & 1234 & 1 & 2 & 0\\
3241 & 12 13 14 23 & 1324 2134 3214 & 3214 & 1243 & 2 & 3 & 1\\
3412 & 13 14 23 24 & 1324 3412 & 3412 & 1234 & 1 & 2 & 0\\
3421 & 12 13 14 23 24 & 1324 2134 3214 3412 & 3412 & 1243 & 2 & 3 & 1\\
4123 & 14 24 34 & 1243 & 1243 & 3124 & 3 & 3 & 2\\
4132 & 14 23 24 34 & 1243 1324 1432 & 1432 & 2134 & 2 & 3 & 1\\
4213 & 12 14 24 34 & 1243 2134 2143 & 2143 & 3124 & 3 & 3 & 2\\
4231 & 12 13 14 24 34 & 1243 2134 2143 4231 & 4231 & 1234 & 1 & 3 & 0\\
4312 & 13 14 23 24 34 & 1243 1324 1432 3412 & 3412 & 2134 & 2 & 3 & 1\\
4321 & all six & all nine & 4321 & 1234 & 1 & 3 & 0\\
\bottomrule
\end{tabular}
\end{table}

The values \(0,1,2,3\) of \(\linv\) occur \(1\), \(9\), \(10\) and \(4\) times. The four elements with \(\linv=3\) are \(2341\), \(2431\), \(4123\) and \(4213\). The maximum is \(3=N-1\). In every row the longest involution prefix is unique, \(\linv(w)\le\disp(w)\), and \(\disp(\rem(w))<\disp(w)\).

Two rows illustrate Lemma~\ref{lem:K} and Proposition~\ref{prop:drop}.
\begin{itemize}
\item \(w=3412\): \(D=2\), with largest excedances at \(p=1\) (\(v=3\)) and \(p=2\) (\(v=4\)). For \(p=1\): \(L=\{3\}\), \(R=\{1,2\}\), \(h(3)=1\), \(h(2)=0\), so \(c_0=2\), \(J=\{2,3\}\) and \(\iota=(2\ 3)=1324\). This is an involution prefix in the table, and it moves \(v=3\). For \(p=2\): \(L=\{3,4\}\), \(R=\{1,2\}\), \(h(4)=1\), \(h(3)=2\), \(h(2)=1\), \(h(1)=0\), so \(c_0=1\), \(J=[1,4]\) and \(\iota=(1\ 3)(2\ 4)=3412=w\).
\item \(w=2431\): \(D=3\), and the only entry of displacement \(3\) is the drop \(w(4)=1\). Then \(\varphi(w)=4213\) has the excedance \(p=1\), \(v=4\): \(L=\{4\}\), \(R=\{1,2,3\}\), \(h(4)=1\), \(h(3)=0\), \(c_0=3\) and \(\iota=(3\ 4)=1243\), an involution prefix of \(4213\) with \(\iota(4)=3\). The table gives \(a(4213)=2143\), and \(\rem(4213)=3124\), whose first entry is \(3\neq4\). Correspondingly, \(\varphi(3124)=1342=\rem(2431)\) has no drop of size \(3\).
\end{itemize}

\subsection{The worked example of Hart and Rowley}
Hart and Rowley give the example \(w=r_6r_3r_2r_1r_4r_5\) in type \(A_6\), with ancestor decomposition \((r_3r_6)(r_2r_4)(r_1r_5)\)~\cite{HR}. We reproduce it with our conventions, taking \(r_i=s_i\in S_7\); this checks the conventions. Multiplying from the right end gives \(w=4125736\): start with \(s_5=1234657\), then exchange the values \(4,5\); then \(1,2\); then \(2,3\); then \(3,4\); then \(6,7\). So
\[
\VInv(w)=V=\{(1,4),(2,4),(3,4),(3,5),(3,7),(6,7)\}.
\]
Let \(\iota\) be an involution with \(\PInv(\iota)\subseteq V\). By Lemma~\ref{lem:invinv}(1), each \(2\)-cycle of \(\iota\) is a pair in \(V\). Lemma~\ref{lem:invinv}(2) excludes four of them.
\begin{itemize}
\item \((1\ 4)\): for \(c=2\) and \(c=3\) the pair \((1,c)\) is not in \(V\), so \(\iota(c)>4\). Then \((2\ \iota(2))\) is a \(2\)-cycle with \(\iota(2)\ge5\), but \(V\) has no pair \((2,b)\) with \(b\ge5\).
\item \((2\ 4)\): \((2,3)\notin V\), so \(\iota(3)>4\), and the \(2\)-cycle \((3\ \iota(3))\) is \((3\ 5)\) or \((3\ 7)\). Then \(\iota(2)=4>3=\iota(\iota(3))\), so \((2,\iota(3))\) is an inversion of \(\iota\). It is \((2,5)\) or \((2,7)\), not in \(V\).
\item \((3\ 5)\): take \(c=4\). If \(\iota(4)>3\), then \((4,5)\) is an inversion, not in \(V\). So \(\iota(4)<3\), and \((\iota(4)\ 4)\) is \((1\ 4)\) or \((2\ 4)\), both excluded.
\item \((3\ 7)\): take \(c=4\). If \(\iota(4)>3\), then \((4,7)\) is an inversion, not in \(V\). So \(\iota(4)<3\), excluded as before.
\end{itemize}
So the \(2\)-cycles of \(\iota\) are among \((3\ 4)\) and \((6\ 7)\). By Lemma~\ref{lem:invinv}(3), \(s_3\), \(s_6\) and \(s_3s_6\) are all involution prefixes. Hence \(a(w)=s_3s_6\) and \(\rem(w)=3125647\).

Next, \(\VInv(3125647)=\{(1,3),(2,3),(4,5),(4,6)\}\). A \(2\)-cycle \((1\ 3)\) is excluded: \((1,2)\) is not in the set, so \(\iota(2)>3\), and no pair \((2,b)\) with \(b\ge4\) is in the set. A \(2\)-cycle \((4\ 6)\) is excluded: if \(\iota(5)>4\), then \((5,6)\) is an inversion, not in the set; so \(\iota(5)<4\), and no pair \((b,5)\) with \(b\le3\) is in the set. So the involution prefixes are \(s_2\), \(s_4\) and \(s_2s_4\), the ancestor is \(s_2s_4\), and the remainder is \(2134657=s_1s_5\), an involution. This is the decomposition \((r_3r_6)(r_2r_4)(r_1r_5)\) of Hart and Rowley, with \(\linv(w)=3\). The displacement takes the values \(3,2,1\) along \(4125736\), \(3125647\), \(2134657\), as Corollary~\ref{cor:A} predicts.

\subsection{\texorpdfstring{\(B_2\)}{B2}}
In the labels \(\lambda\), \(B_2\) consists of the eight \(w\in S_4\) with \(w(5-i)=5-w(i)\): \(1234\), \(1324=s_0\), \(2143=s_1\), \(3142=s_0s_1\), \(2413=s_1s_0\), \(3412=s_0s_1s_0\), \(4231=s_1s_0s_1\) and \(4321\). The two elements of length \(2\) have the single involution prefix \(s_0\), respectively \(s_1\), and involution length \(2\). The other five non-identity elements are involutions. Table~\ref{tab:S4} gives the same in \(S_4\): \(a(3142)=1324\) with remainder \(2143\); \(a(2413)=2143\) with remainder \(1324\); and \(1324\), \(2143\), \(3412\), \(4231\), \(4321\) are involutions. So \(\linv^B=\linv^A\) on \(B_2\), as Theorem~\ref{thm:B} says.

\section{Computer corroboration}\label{sec:computer}

\textbf{No step of any proof in this paper uses a computation.} The computations below were made to catch errors. They are corroboration only. They come from a single implementation, so they are not certificates in the stronger sense of two independent programs that agree.

\emph{Programs and environment.} Two Python scripts, run with Python 3.12.10 and the standard library only, on 2026-11-27. All arithmetic is exact integer arithmetic. The scripts and their archived outputs are in the directory \file{scripts/} next to this paper, in the folder \file{typeA/} of the repository \url{https://github.com/nebulousice/hart-rowley-conjecture-1-3}:
\begin{itemize}
\item \file{scripts/small_cases.py}, with outputs \file{scripts/out_small_cases_N6.txt} and \file{scripts/out_small_cases_N7.txt};
\item \file{scripts/typeB_cases.py}, with output \file{scripts/out_typeB_cases_n4.txt}.
\end{itemize}

\emph{Commands} (run in the directory \file{scripts/} next to this paper; see also \file{REPRODUCE.md}):
\begin{verbatim}
python small_cases.py 6 tables > out_small_cases_N6.txt   # 11 s
python small_cases.py 7        > out_small_cases_N7.txt   # not timed
python typeB_cases.py 4        > out_typeB_cases_n4.txt   # 16 s
\end{verbatim}

\emph{What is checked, type A.} For every \(w\in S_N\) with \(2\le N\le7\), the script computes all involution prefixes from the definition (length additivity against every involution) and the ancestor decomposition. It checks that the longest involution prefix is unique, that \(\disp(\rem(w))<\disp(w)\) for \(w\neq e\), that \(\linv(w)\le\disp(w)\) and \(\linv(w)\le N-1\), and that \(\linv(s_1\cdots s_{N-1})=N-1\). For \(w\neq e\) it checks Lemma~\ref{lem:K} at every pair \((w,p)\) with \(w(p)=p+\disp(w)\), and at every pair \((\varphi(w),N+1-p)\) with \(w(p)=p-\disp(w)\) (one test for each position of largest displacement). It also checks the instances of Theorem~\ref{thm:GG}(2) and Lemma~\ref{lem:compl}. For \(N\le5\) it checks Lemmas~\ref{lem:product} and~\ref{lem:prefixcrit} on all pairs. For \(N\le6\) it builds the Bruhat order from Definition~\ref{def:bruhat} and checks Lemma~\ref{lem:rankmatrix}. For \(N=7\) it compares elements in Bruhat order by the rank-matrix criterion in both directions, that is, it also uses the converse of Lemma~\ref{lem:rankmatrix}~\cite[Ch.~2]{BB}; this converse is used only in the computation, never in a proof. It also prints Tables~\ref{tab:S3} and~\ref{tab:S4}, which agree with the hand computation. Every failure counter in both output files is \(0\). Table~\ref{tab:hist} records the distribution of \(\linv\) and the number of tests of Lemma~\ref{lem:K}.

\begin{table}[ht]
\centering
\small
\caption{Archived output \texttt{out\_small\_cases\_N7.txt}: number of \(w\in S_N\) with \(\linv(w)=0,1,\dots\), and the number of tests of Lemma~\ref{lem:K} (one per position \(p\) with \(\abs{w(p)-p}=\disp(w)\), over all \(w\neq e\); key \texttt{lemmaK\_pairs}).}\label{tab:hist}
\begin{tabular}{rrlrr}
\toprule
\(N\) & \(\abs{S_N}\) & distribution of \(\linv\) & \(\max\linv\) & tests of Lemma~\ref{lem:K}\\
\midrule
2 & 2 & 1, 1 & 1 & 2\\
3 & 6 & 1, 3, 2 & 2 & 8\\
4 & 24 & 1, 9, 10, 4 & 3 & 38\\
5 & 120 & 1, 25, 48, 36, 10 & 4 & 182\\
6 & 720 & 1, 75, 240, 264, 104, 36 & 5 & 1064\\
7 & 5040 & 1, 231, 1212, 1906, 1098, 438, 154 & 6 & 7176\\
\bottomrule
\end{tabular}
\end{table}

\emph{What is checked, types B and D.} For \(n\le4\) the script computes \(\ell_B\) and \(\ell_D\) as word lengths by breadth-first search, and ancestors in \(S_{2n}\) over all involutions of \(S_{2n}\). It checks Lemma~\ref{lem:lAlB}, the agreement of the weak orders (Lemma~\ref{lem:Bweak}), \(a_B=a_A\) (Proposition~\ref{prop:Bancestor}) and \(\linv^B=\linv^A\le\disp_A\) (Theorem~\ref{thm:B}); all failure counters are \(0\). The distributions of \(\linv^B\) are \(1,1\) for \(B_1\); \(1,5,2\) for \(B_2\); \(1,19,20,8\) for \(B_3\); and \(1,75,156,123,29\) for \(B_4\). For type D it finds \(32\) pairs \((u,w)\) in \(D_3\) and \(788\) in \(D_4\) with \(u\leR^D w\) but not \(u\leR^A w\); \(6\) ancestor mismatches and \(4\) involution-length mismatches in \(D_3\); \(55\) and \(33\) in \(D_4\); and, as the shortest mismatch in \(D_3\), the element with window \((-3,-1,2)\) of Remark~\ref{rem:D}.

\section{Formal verification}\label{sec:lean}

Nothing in this paper is formalized in Lean or in any other proof assistant. In particular, Theorems~\ref{thm:GG}, \ref{thm:A} and~\ref{thm:B} are not formalized. The Bruhat order of a general Coxeter group, which Appendix~\ref{app:GG} needs, is not available in Mathlib as far as we know. The only Lean formalization connected with Hart--Rowley Conjecture~1.3 in this repository is that of the \(H_3\) counterexample, in the companion note (files under \file{../lean/}, in the parent directory of this paper).

\section{Status of this work}\label{sec:status}

Disclosure: Human researcher assisted by AI systems during search, proof drafting, and verification. Internally audited and computationally verified, externally unrefereed.

Here the computations are corroboration only; no proof step uses them (Section~\ref{sec:computer}).

\textbf{Expert verification required.} Verified so far with:
\begin{itemize}
\item \emph{Lean 4 and Mathlib:} nothing in this paper (Section~\ref{sec:lean}).
\item \emph{Independent certificates:} none; no proof step uses a computation. A single exact implementation (two Python scripts) corroborates the results for \(S_N\), \(N\le7\), and for \(B_n\), \(D_n\), \(n\le4\) (Section~\ref{sec:computer}).
\item \emph{AI referees}, reading the full proof:
round~1 (commit \texttt{4101231}), Claude Opus: minor issues (five), all fixed; Claude Fable: accept, no issues. In round~1 the theorem of Gaetz and Gao was still cited, not proved.
Round~2 (commit \texttt{599cc84}), after Appendix~\ref{app:GG} was added: Claude Opus, Claude Fable and Claude Sonnet each accept with no issues.
The commit hashes refer to the private working repository in which the review was done.
\end{itemize}

\appendix

\section{Proof of the theorem of Gaetz and Gao}\label{app:GG}

\emph{Credit.} Everything in this appendix is the argument of Gaetz and Gao~\cite[\S3]{GG}: their Lemma~3.1, Theorem~1.1 and Corollary~1.2. We write it out in our notation and our own words, and we justify in full each step that the source states briefly. The first paragraph of their proof of Lemma~3.1 is Lemma~\ref{lem:compl} above. There is one change of substance: their proof of Theorem~1.1 uses the finiteness of Bruhat intervals; we replace this by a length argument (Theorem~\ref{thm:lr}).

\emph{Inputs.} \((W,S)\) is an arbitrary Coxeter system. We use only Facts~\ref{fact:length}, \ref{fact:lifting} and~\ref{fact:subword}, the bound~\eqref{eq:I2}, and Lemmas~\ref{lem:inverse}, \ref{lem:lift}, \ref{lem:suffix} and~\ref{lem:compl}. None of these depends on Theorem~\ref{thm:GG}. For \(x,y\in W\) put
\begin{gather*}
\LB(x,y)=\{g\mid g\leL x,\ g\leB y\},\qquad
\BR(x,y)=\{g\mid g\leB x,\ g\leR y\},\\
\LR(x,y)=\{g\mid g\leL x,\ g\leR y\}.
\end{gather*}
Each contains \(e\). A \emph{greatest element} of a set is an element of the set that lies \(\geB\) every element of it. It is unique if it exists, because \(\leB\) is a partial order.

\subsection{The downward Demazure operator}
For \(s\in S\) and \(g\in W\) let \(\delta_s(g)=gs\) if \(\ell(gs)<\ell(g)\), and \(\delta_s(g)=g\) otherwise.

\begin{lemma}\label{lem:demazure}
For \(s\in S\) and \(g,h\in W\):
\begin{enumerate}
\item \(\delta_s(g)\leB g\);
\item if \(g\leB h\), then \(\delta_s(g)\leB\delta_s(h)\).
\end{enumerate}
\end{lemma}

\begin{proof}
(1) If \(\delta_s(g)=gs\) with \(\ell(gs)<\ell(g)\), then \(g=(gs)s\) is a chain step, so \(gs\leB g\). Otherwise \(\delta_s(g)=g\).

(2) If \(\ell(hs)>\ell(h)\), then \(\delta_s(h)=h\), and \(\delta_s(g)\leB g\leB h\) by (1). Let \(\ell(hs)<\ell(h)\), so \(\delta_s(h)=hs\). If \(\ell(gs)<\ell(g)\), then \(\delta_s(g)=gs\), and (R1) of Lemma~\ref{lem:lift} gives \(gs\leB hs\). If \(\ell(gs)>\ell(g)\), then \(\delta_s(g)=g\), and (R2) gives \(g\leB hs\).
\end{proof}

\subsection{Sweeps}
Fix a word \(s_1s_2\cdots s_m\) in the generators. For a set of indices \(I=\{i_1<\dots<i_q\}\subseteq\{1,\dots,m\}\) and \(g\in W\), the \emph{sweep} of \(g\) along \(I\) is
\[
\Delta_I(g)=\delta_{s_{i_1}}\bigl(\delta_{s_{i_2}}(\cdots\delta_{s_{i_q}}(g)\cdots)\bigr).
\]
The operators are applied in decreasing order of index, the one with the largest index first. Also \(\Delta_\emptyset(g)=g\).

\begin{lemma}\label{lem:sweep}
\begin{enumerate}
\item Each \(\Delta_I\) is monotone for \(\leB\), and \(\Delta_I(g)\leB g\).
\item If \(I\subseteq J\), then \(\Delta_J(g)\leB\Delta_I(g)\) for every \(g\).
\item Let \(w\leL x\), and let \(j_1<\dots<j_q\) be indices such that \(s_{j_1}\cdots s_{j_q}\) is a reduced word for \(w\). Then \(\Delta_{\{j_1,\dots,j_q\}}(x)=xw^{-1}\), and each of the \(q\) operators lowers the length by exactly one.
\end{enumerate}
\end{lemma}

\begin{proof}
(1) A composition of monotone maps is monotone. A composition of maps with \(\delta_s(g)\leB g\) has the same property. Both follow from Lemma~\ref{lem:demazure} and transitivity.

(2) By induction on \(\abs{J\setminus I}\), it suffices to treat \(J=I\cup\{j\}\) with \(j\notin I\). Let \(I_>\) and \(I_<\) be the elements of \(I\) larger and smaller than \(j\). Then \(\Delta_J=\Delta_{I_<}\circ\delta_{s_j}\circ\Delta_{I_>}\) and \(\Delta_I=\Delta_{I_<}\circ\Delta_{I_>}\). Put \(h=\Delta_{I_>}(g)\). By Lemma~\ref{lem:demazure}(1), \(\delta_{s_j}(h)\leB h\). By (1) applied to \(\Delta_{I_<}\), \(\Delta_{I_<}(\delta_{s_j}(h))\leB\Delta_{I_<}(h)\), which is the claim.

(3) For \(0\le p\le q\) let \(w_p=s_{j_1}\cdots s_{j_p}\), so \(w_0=e\) and \(w_q=w\). Let \(g_p\) be the element obtained from \(x\) by applying \(\delta_{s_{j_q}},\dots,\delta_{s_{j_{p+1}}}\), so \(g_q=x\). We show by downward induction on \(p\) that
\[
g_p=xw^{-1}w_p\qquad\text{and}\qquad\ell(g_p)=\ell(xw^{-1})+p.
\]
For \(p=q\): \(g_q=x=xw^{-1}w\), and \(\ell(x)=\ell(xw^{-1})+\ell(w)\) because \(w\leL x\). Let \(1\le p\le q\) and assume the statement for \(p\). Then \(g_ps_{j_p}=xw^{-1}w_{p-1}\). By Fact~\ref{fact:length}, \(\ell(g_ps_{j_p})\le\ell(xw^{-1})+\ell(w_{p-1})\le\ell(xw^{-1})+p-1<\ell(g_p)\). So \(s_{j_p}\) is a right descent of \(g_p\), \(g_{p-1}=g_ps_{j_p}=xw^{-1}w_{p-1}\), and \(\ell(g_{p-1})=\ell(g_p)-1\) by Fact~\ref{fact:length}. For \(p=0\) this gives \(g_0=xw^{-1}\).
\end{proof}

\subsection{Greatest elements of \texorpdfstring{\(\LB\)}{LB} and \texorpdfstring{\(\BR\)}{BR}}

\begin{lemma}[Gaetz--Gao, Lemma~3.1]\label{lem:lb}
For all \(x,y\in W\), the set \(\LB(x,y)\) has a greatest element \(\lb(x,y)\), and \(\BR(x,y)\) has a greatest element \(\br(x,y)=\lb(y^{-1},x^{-1})^{-1}\).
\end{lemma}

\begin{proof}
We construct a candidate \(z\) by sweeping \(x\) along a reduced word of \(y\), show that \(z\in\LB(x,y)\), and then use complementation to show that it is greatest. Fix a reduced word \(s_1\cdots s_m\) for \(y\), with \(m=\ell(y)\), and take sweeps along subsets of \(\{1,\dots,m\}\).

\emph{Construction.} Put \(u_m=x\) and \(u_{i-1}=\delta_{s_i}(u_i)\) for \(i=m,m-1,\dots,1\). So \(u_0=\Delta_{\{1,\dots,m\}}(x)\). Let \(i_1<\dots<i_r\) be the indices \(i\) with \(u_{i-1}\neq u_i\). At each of them \(u_{i-1}=u_is_i\) and \(\ell(u_{i-1})=\ell(u_i)-1\); at every other index \(u_{i-1}=u_i\). Since the operators act from the largest index down, \(u_0=xs_{i_r}s_{i_{r-1}}\cdots s_{i_1}\). Put
\[
z=s_{i_1}s_{i_2}\cdots s_{i_r},\qquad\text{so that}\qquad u_0=xz^{-1}\quad\text{and}\quad\ell(u_0)=\ell(x)-r.
\]

\emph{\(z\in\LB(x,y)\).} From \(x=u_0z\) and Fact~\ref{fact:length}, \(\ell(x)\le\ell(u_0)+\ell(z)\le(\ell(x)-r)+r=\ell(x)\). Both inequalities are equalities. So \(\ell(z)=r\), that is, the word \(s_{i_1}\cdots s_{i_r}\) is reduced; and \(\ell(x)=\ell(xz^{-1})+\ell(z)\), that is, \(z\leL x\). This reduced word for \(z\) is a subword of the reduced word \(s_1\cdots s_m\) for \(y\). By Fact~\ref{fact:subword}, \(z\leB y\).

\emph{\(z\) is greatest.} Let \(w\in\LB(x,y)\). Since \(w\leB y\), Fact~\ref{fact:subword} gives indices \(j_1<\dots<j_q\) such that \(s_{j_1}\cdots s_{j_q}\) is a reduced word for \(w\). Put \(I=\{j_1,\dots,j_q\}\). Since \(w\leL x\), Lemma~\ref{lem:sweep}(3) gives \(\Delta_I(x)=xw^{-1}\). Lemma~\ref{lem:sweep}(2), with \(I\subseteq\{1,\dots,m\}\), gives
\[
u_0=\Delta_{\{1,\dots,m\}}(x)\leB\Delta_I(x)=xw^{-1}.
\]
Both \(u_0\) and \(xw^{-1}\) are prefixes of \(x\): \(\ell(x)=\ell(u_0)+\ell(u_0^{-1}x)\) because \(u_0^{-1}x=z\) and \(\ell(x)=\ell(u_0)+\ell(z)\); and \(\ell(x)=\ell(xw^{-1})+\ell(w)\) with \((xw^{-1})^{-1}x=w\). So Lemma~\ref{lem:compl}, applied to the triple \((u_0,xw^{-1},x)\), gives \((xw^{-1})^{-1}x\leB u_0^{-1}x\), that is, \(w\leB z\). So \(z\) is the greatest element of \(\LB(x,y)\).

\emph{The set \(\BR\).} By Lemma~\ref{lem:inverse}, \(g\leB x\) if and only if \(g^{-1}\leB x^{-1}\), and \(g\leR y\) if and only if \(g^{-1}\leL y^{-1}\). So \(g\mapsto g^{-1}\) maps \(\BR(x,y)\) bijectively onto \(\LB(y^{-1},x^{-1})\), and it preserves \(\leB\). Hence \(\lb(y^{-1},x^{-1})^{-1}\) is the greatest element of \(\BR(x,y)\).
\end{proof}

\subsection{The greatest element of \texorpdfstring{\(\LR\)}{LR}}

\begin{theorem}[Gaetz--Gao, Theorem~1.1]\label{thm:lr}
For all \(x,y\in W\), the set \(\LR(x,y)\) has a greatest element \(\lr(x,y)\). Moreover \(\lr(y^{-1},x^{-1})=\lr(x,y)^{-1}\).
\end{theorem}

\begin{proof}
We alternate the two operations of Lemma~\ref{lem:lb} until the result stabilizes. Define \(z_0=x\) and, for \(i\ge0\),
\[
z_{2i+1}=\br(z_{2i},y),\qquad z_{2i+2}=\lb(x,z_{2i+1}).
\]
These exist by Lemma~\ref{lem:lb}.

\emph{The sequence decreases and becomes constant.} \(z_{2i+1}\in\BR(z_{2i},y)\), so \(z_{2i+1}\leB z_{2i}\). Also \(z_{2i+2}\in\LB(x,z_{2i+1})\), so \(z_{2i+2}\leB z_{2i+1}\). Hence \(z_{j+1}\leB z_j\) for every \(j\). By~\eqref{eq:I2}, \(\ell(z_{j+1})\le\ell(z_j)\), with equality only if \(z_{j+1}=z_j\). A non-increasing sequence of non-negative integers is eventually constant. So there is \(j_0\) with \(z_{j+1}=z_j\) for all \(j\ge j_0\); call the common value \(z\).

\emph{\(z\in\LR(x,y)\).} Choose \(i\) with \(2i\ge j_0\). Then \(z=z_{2i+1}=\br(z_{2i},y)\in\BR(z,y)\), so \(z\leR y\). And \(z=z_{2i+2}=\lb(x,z_{2i+1})\in\LB(x,z)\), so \(z\leL x\).

\emph{\(z\) is greatest.} Let \(g\in\LR(x,y)\). We show \(g\leB z_j\) for all \(j\), by induction on \(j\). For \(j=0\): \(g\leL x\), so \(g\leB x=z_0\) by Lemma~\ref{lem:suffix}(1). If \(g\leB z_{2i}\), then, since \(g\leR y\), \(g\in\BR(z_{2i},y)\); so \(g\leB\br(z_{2i},y)=z_{2i+1}\). If \(g\leB z_{2i+1}\), then, since \(g\leL x\), \(g\in\LB(x,z_{2i+1})\); so \(g\leB\lb(x,z_{2i+1})=z_{2i+2}\). Hence \(g\leB z\).

\emph{Symmetry.} By Lemma~\ref{lem:inverse}(1), \(g\leL x\) if and only if \(g^{-1}\leR x^{-1}\), and \(g\leR y\) if and only if \(g^{-1}\leL y^{-1}\). So inversion maps \(\LR(x,y)\) bijectively onto \(\LR(y^{-1},x^{-1})\). By Lemma~\ref{lem:inverse}(2) it preserves \(\leB\), so it maps the greatest element to the greatest element.
\end{proof}

\subsection{Proof of Theorem~\ref{thm:GG}}

\begin{proof}[Proof of Theorem~\ref{thm:GG} (Gaetz--Gao, Corollary~1.2)]
Let \(w\neq e\) and put \(a=\lr(w^{-1},w)\), the greatest element of
\[
\X=\LR(w^{-1},w)=\{g\mid g\leL w^{-1},\ g\leR w\}.
\]
\begin{itemize}
\item \emph{\(a=a^{-1}\).} Theorem~\ref{thm:lr} with \((x,y)=(w^{-1},w)\) has \((y^{-1},x^{-1})=(w^{-1},w)\). So \(\lr(w^{-1},w)=\lr(w^{-1},w)^{-1}\).
\item \emph{\(a\neq e\).} Let \(s\) be the first letter of a reduced word for \(w\). Then \(s\leR w\), and by Lemma~\ref{lem:inverse}(1), \(s=s^{-1}\leL w^{-1}\). So \(s\in\X\), and \(s\leB a\) by Theorem~\ref{thm:lr}. By~\eqref{eq:I2}, \(\ell(a)\ge\ell(s)=1\).
\item So \(a^2=e\) and \(a\neq e\): \(a\) is an involution. Also \(a\in\X\), so \(a\leR w\).
\item \emph{Domination.} Let \(\iota\) be an involution with \(\iota\leR w\). By Lemma~\ref{lem:inverse}(1), \(\iota=\iota^{-1}\leL w^{-1}\). So \(\iota\in\X\), and \(\iota\leB a\) by Theorem~\ref{thm:lr}.
\item \emph{Uniqueness.} For such \(\iota\), \eqref{eq:I2} gives \(\ell(\iota)\le\ell(a)\), with equality only if \(\iota=a\). So \(a\) is the unique involution of largest length with \(a\leR w\).
\end{itemize}
The last two items are statements (2) and (1) of Theorem~\ref{thm:GG}, with \(a(w)=a\).
\end{proof}

For example, in \(S_3\) with \(w=231\): the suffixes of \(w^{-1}=312=s_2s_1\) are \(e\), \(s_1\) and \(312\), and the prefixes of \(231=s_1s_2\) are \(e\), \(s_1\) and \(231\), since each of these elements has only one reduced word. So \(\X=\{e,213\}\) and \(a(231)=213\), as in Table~\ref{tab:S3}.

\begin{thebibliography}{9}

\bibitem{BB}
A.~Bj\"orner and F.~Brenti, \emph{Combinatorics of {C}oxeter groups}, Graduate Texts in Mathematics, vol.~231, Springer, New York, 2005, \href{https://doi.org/10.1007/3-540-27596-7}{doi:10.1007/3-540-27596-7}.

\bibitem{GG}
C.~Gaetz and Y.~Gao, \emph{Separable elements and {B}ruhat orders}, preprint, 2026, \href{https://arxiv.org/abs/2610.22857}{arXiv:2610.22857}.

\bibitem{HR}
S.~B. Hart and P.~J. Rowley, \emph{A note on involution prefixes in {C}oxeter groups}, Arch. Math. (Basel) \textbf{126} (2026), no.~3, 231--237, \href{https://doi.org/10.1007/s00013-025-02220-z}{doi:10.1007/s00013-025-02220-z}; preprint \href{https://arxiv.org/abs/2502.00777}{arXiv:2502.00777}.

\bibitem{Lev}
Salem Levinstein, \emph{Counterexamples to the {H}art--{R}owley {I}nvolution-{L}ength {C}onjecture in {F}inite {C}oxeter {G}roups}, Zenodo preprint, 2026-11-27, \href{https://doi.org/10.5281/zenodo.24581942}{doi:10.5281/zenodo.24581942}.

\end{thebibliography}

\end{document}
